% Modernised reading — fonds Alexandre Grothendieck, shelfmark 82
% (pages 1-78, the whole folder).
%
% This is an INTERPRETATION, not a transcription. It re-reads the pages in
% today's notation and vocabulary, states the hypotheses the manuscript leaves
% implicit, and drops the critical apparatus so that the mathematics reads
% straight through. Everything the brackets and underlines carried moves into a
% footnote. It is held to being mathematically correct as it stands; where the
% manuscript is loose or mistaken, the true statement is what appears here and
% the footnote says what the page has. In any dispute of substance the
% transcription governs: whoever cites Grothendieck must cite batch-01.fr.tex
% (pp. 1-20), batch-02.fr.tex (pp. 21-40), batch-03.fr.tex (pp. 41-60) or
% batch-04.fr.tex (pp. 61-78).
%
% Pass: Opus 5.5 (claude-opus-5-5), 2026-10-10 — first pass, unchecked against the pages by a human.
% The folder was read whole, its four batches together; this is one document
% for the shelfmark. It was made on Opus 5.5 and has not been measured against
% the reference reading of folder 115 (made on Opus 5). The literature named
% in the footnotes is cited from memory and was not checked against the
% sources. The computations of pp. 2-25 (Jacobi identity, Cramer, Casimir
% forms, the example) were re-done symbolically for this reading.
%
% Scope. Three runs, each opened by a pink cover: p. 1 « Algèbres de Lie de
% rang 3 associées aux formes bil. symétriques de rang 3 », « (29 p) »
% (pp. 2-31); p. 32, no title, « (21 p) » (pp. 33-55, i.e. 21 written pages
% once the blanks 35 and 45 are set aside); p. 56 « Dictionnaire lié à
% GP(1) ≃ SO(3) » (pp. 57-78). Pages 35 and 45 are blank and not described.
%
% Spine. The exceptional isomorphism PGL_2 ≃ SO(3), over an arbitrary base
% scheme, Spec Z included, and in all its twisted forms:
%   pp. 2-27   a bilinear form phi on a rank-3 module E with a basis omega of
%              det E gives a bracket on the dual E'; symmetric phi gives Lie
%              algebras; f : E -> E', its cofactor map f' = Lambda^2 f, the
%              bracket on E, the action of E' on E; the example xy + z^2
%              gives sl(2) -> gp(1); Casimir forms; rescaling; char. 2.
%   pp. 28-31  G = SO(q), its simply connected cover G~ (kernel mu_2), and
%              E' ≃ Lie G proved; E ≃ Lie G~ announced, « il faut ouvrir une
%              parenthèse ».
%   pp. 33-55  second notebook: the rotation of a regular n-gon and F_n;
%              special subalgebras of a quaternion algebra = lines of the
%              plane = degree-2 divisors of the conic = smooth commutative
%              one-dimensional subgroups; centralisers (pp. 36-43, char. 2
%              counterexamples); X ≃ P(V), the quadric of rank-one
%              endomorphisms ≃ X x X, the canonical involution, the
%              « jacobian » subgroups (pp. 46-55).
%   pp. 57-78  the Dictionnaire: nine fibred categories (gerbes), typical
%              objects over Z, actions of G_0 = PGL_2, the functors between
%              them, then the quaternion algebra M as the key: G, G~, g, g~
%              from M, the trace form, q = -det on g~, omega, f and f',
%              Borel subgroups, and in char. 2 the Frobenius X -> X^(2).
%
% Notation fixed for the whole folder. S a scheme; E, E' locally free of rank
% 3 in duality; omega a basis of det E, omega' the dual basis of det E';
% alpha : Lambda^2 E ≃ E', alpha' : Lambda^2 E' ≃ E (his (6)-(7)); phi a
% bilinear form on E, f : E -> E' with <y, f(x)> = phi(x, y); [ , ]_phi the
% bracket on E', f' = alpha' o Lambda^2 f o alpha^{-1} : E' -> E, phi' its
% form, [ , ]_phi' the bracket on E; Delta the scalar with delta(phi) = Delta
% omega^{-2}. rho_E' = ad on E', rho_E its contragredient on E, rho~_E = ad
% on E. q a quadratic form with phi(x,x) = 2q(x); (q, omega) special
% orthogonal when delta'(q) = -omega^{-2}. G = SO(q), G~ its simply connected
% cover, g = Lie G, g~ = Lie G~. P-check(F) = line subbundles of F, P(F) =
% hyperplanes of F (his convention, kept). M a quaternion algebra, M^x its
% unit group (his V(M)^* and round L(M)), Nrd = det, Trd = Tr. X the conic,
% P the plane, Y a degree-2 divisor of X, A a special subalgebra (rank 2,
% containing 1), T = A^x/G_m the jacobian subgroup, h = Lie T. The pages'
% H for a one-dimensional subgroup (pp. 36-43) is written T here, since it
% is the same object and since H is also a basis vector of gp(1). The form
% -det on M is Q_M, the quadric of rank-one endomorphisms is the calligraphic
% Q. X, Y, H also name the basis of gp(1) in section I, as on the pages; the
% conic X and the divisor Y appear only from p. 34 on.
%
% Substantive departures from the pages, with page numbers:
%  - p. 6, (23 bis): c = 1 should be c = 2 (the matrix beside it has 2).
%  - p. 9, (38): f' = alpha' o Lambda^2 f o alpha^{-1}; the page has alpha.
%  - p. 10: the brackets' attributions « sur E », « sur E' » are swapped
%    relative to (12) and (40 bis).
%  - p. 13: the example (21) has Delta = -2, not 2 (p. 18 has -2).
%  - pp. 15, 17: the bracket on E' is [ , ]_phi (12); the subscript in
%    (55), (57), (59 bis) is phi' or unclear. (52) => (53) also needs
%    Tr rho_E = 0, the (60) of the margin of p. 18; supplied.
%  - p. 19, (66)-(69): sign. f'(X) = -2X~, f'(Y) = -2Y~, f'(H) = -H~, so
%    phi' and q' are the negatives of the page's. Confirmed by (54)-(55),
%    p. 75, ff' = -2. Consequences: p. 21 « f' = 2 id » is -2 id; p. 22
%    (74) first line is Cas_E'(x', y') = -2 phi'(x', y'), and (77) likewise;
%    p. 23 phi'(H,H) = -1. The second line of (74), (75), (76) are right.
%  - p. 21: « groupe dérivé » of sl(2) is the derived algebra.
%  - p. 25: in the proof of (80) det f' = det g' gives Delta(f)^2 =
%    Delta(g)^2; (80) and (81) need Delta invertible; in (81) lambda must be
%    the inverse of a square.
%  - p. 37: « g = GP(V) » is g = Lie GP(V). p. 39: « alpha lambda ≡ 1 »
%    read lambda alpha = alpha.
%  - p. 42, b): needs p invertible on S; at a fibre of characteristic p,
%    N/H is all of F_p^x, which is {±1} only for p = 3.
%  - p. 46: « rg 2 » is rank 3. p. 48: g-bar is g~. p. 49: Gamma(u) is the
%    form of (5).
%  - p. 55: « T~ is the centraliser of h » holds when 2 is invertible; in
%    char. 2 the centraliser is bigger (p. 40).
%  - p. 64: the stabiliser (11) is the transpose of (10); s_0 and « ∞ » name
%    the same point; footnoted, conventions explained.
%  - pp. 66: the « ∼ » of (23)-(24) are equivalences over Z[1/2] only.
%  - p. 71: the pairing induced by Phi_1 is equal to (41), not only up to
%    sign. p. 72, (47): det(V-check ⊗ V) ≃ (det V-check)^2 ⊗ (det V)^2.
%    p. 73, (48)-(49): omega^{⊗2} read omega^{⊗(-2)} as in (33), (82).
%  - p. 76, (57): the Borel of M^x is the unit group of the rank-3
%    subalgebra L_1; the page writes L~.
%
% Announced and never established in the folder: the construction of G~ and
% the isomorphism E ≃ g~ with its properties for a general special quadratic
% module (p. 31, « ouvrir une parenthèse ») — in substance supplied by
% section VI (pp. 68-75) for E = g~, with the sign of (52) « sans
% garantie »; the equivalences M |-> G, G~, (E, q, omega), (P, X), X (p. 46)
% and the isomorphisms (7) G_0 ≃ Aut(X_0) (p. 62, « nous admettrons »); the
% 72 equivalences and 504 transitivity isomorphisms (p. 58); the full proof
% of the proposition on subgroups of order p (pp. 41-43); whether every
% smooth one-dimensional subgroup is commutative over a base, and the
% normaliser (pp. 43-44, breaks off); the question in the margin of p. 77.
% The folder ends at p. 78 in the middle of section VI.
\documentclass[11pt,a4paper]{article}
\input{../preamble/grothendieck}

\folder{82}
\pages{1}{78}
\dating{[vers 1982]}
\watermark{Édition de démonstration\\interprétation personnelle de l'œuvre}
\foldertitle{SO(3) ≃ GP(1) : notes manuscrites (s.d.) — lecture modernisée du dossier entier}

\begin{document}

\section*{Résumé}

\begin{resume}
Ces pages traitent d'une coïncidence : deux groupes qui n'ont, à première vue,
rien à voir sont le même. L'un est le groupe des rotations de l'espace à trois
dimensions, plus précisément le groupe des transformations qui préservent une
forme quadratique en trois variables, comme $xy + z^2$ ; on le note
$\mathrm{SO}(3)$. L'autre est le groupe des transformations d'une droite
projective, celles qu'on écrit $t \mapsto (at + b)/(ct + d)$ ; Grothendieck le
note $\mathrm{GP}(1)$, on l'écrit aujourd'hui $\mathrm{PGL}_2$. Sur les nombres
complexes la coïncidence est connue depuis longtemps ; l'enjeu ici est de la
rendre vraie partout à la fois — sur les entiers, en toute caractéristique,
caractéristique $2$ comprise, et sur une base quelconque, où les objets peuvent
être « tordus », c'est-à-dire localement standard sans l'être globalement.

L'image classique est celle des quaternions de Hamilton. Un quaternion de norme
$1$ agit sur l'espace des quaternions purs, qui est de dimension $3$, par une
rotation, et deux quaternions opposés donnent la même rotation. Les rotations
sont ainsi des « moitiés » de quaternions, et l'ambiguïté du signe est un
groupe à deux éléments, que Grothendieck note $\mu_2$. Tout le dossier est une
version de cette image où l'on ne divise jamais par $2$ sans le dire.

Le premier cahier part de très peu : un espace de dimension $3$, une forme
bilinéaire, et le choix d'un élément de volume. Avec cela on fabrique un
crochet — un produit du genre du produit vectoriel — sur l'espace dual, et
Grothendieck montre que ce crochet vérifie l'identité de Jacobi dès que la
forme est symétrique. Pour la forme $xy + z^2$, il trouve d'un côté l'algèbre
de Lie de $\mathrm{PGL}_2$, de l'autre celle de $\mathrm{SL}_2$, et entre elles
l'application qui, sur les groupes, est le passage de $\mathrm{SL}_2$ à son
quotient par $\pm 1$. En caractéristique $2$, cette application cesse d'être un
isomorphisme, et les pages en suivent la dégénérescence avec soin.

Le deuxième cahier et le « Dictionnaire » qui clôt le dossier donnent la
version tordue. Neuf sortes d'objets — les formes de $\mathrm{PGL}_2$ et de
$\mathrm{SL}_2$, leurs algèbres de Lie, les algèbres de quaternions, les
coniques, les formes quadratiques de rang $3$ munies d'un volume, les torseurs —
se correspondent toutes, et la clef de toutes les correspondances est
l'algèbre de quaternions : son groupe des unités modulo les scalaires,
ses éléments de norme $1$, ses éléments de trace nulle, la conique de ses
éléments de carré nul. Les deux algèbres de Lie du premier cahier y
reparaissent comme les éléments de trace nulle et les éléments modulo les
scalaires, et la question laissée en suspens à la page 31 y trouve sa réponse.

Il y a aussi un geste qui dit quelque chose de sa manière. Ayant aligné ses
neuf catégories, il compte ce qu'un dictionnaire complet devrait contenir —
$72$ équivalences, $504$ isomorphismes de compatibilité — et écrit : « On fera un
peu moins, mais on essaiera de donner une image claire de la situation. » Le
but n'est pas l'exhaustivité ; c'est que la situation devienne visible.

Les mots à chercher ensuite : isomorphisme exceptionnel
$\mathrm{PGL}_2 \simeq \mathrm{SO}_3$, algèbre de Clifford paire d'une forme
ternaire, algèbres d'Azumaya de rang $4$, gerbes et formes tordues, schéma des
sous-groupes de Borel.

\keywords{exceptional isomorphism PGL2 = SO3, ternary quadratic form, half-discriminant, adjugate matrix, Bianchi classification, Casimir element, Killing form, Heisenberg Lie algebra, affine symplectic group, Spin group, quaternion algebra, Azumaya algebra, even Clifford algebra, reduced norm, canonical involution, Severi-Brauer conic, generalized Jacobian, centralizer of a torus, twisted form, gerbe, nonabelian cohomology, scheme of Borel subgroups, relative Frobenius}
\end{resume}

\pagerange{2}{78}
\section*{Le fil du dossier, et les conventions}

Le dossier est fait de trois suites, ouvertes chacune par une couverture rose,
et il n'y a de sa main aucune pagination. Les nombres au crayon des couvertures
sont des comptes de pages : « (29 p) » sur la première couvre, à une page près,
les pages 2 à 31 ; « (21 p) » sur la deuxième (page 32) tombe juste pour les
pages 33 à 55 une fois écartées les deux pages blanches 35 et 45. Les numéros de
formules sont repris tels que les pages les donnent ; ils forment trois suites :
(1) à (93) aux pages 2 à 29, (1) à (25) aux pages 36 et 46 à 54 (la page 36
ouvre sa propre suite (1)--(2)), (1) à (61) aux pages 60 à 78.

\subsection*{Les stations}

\begin{itemize}
  \item \textbf{I. Algèbres de Lie de rang $3$ (pages 2 à 31).} Une forme
    bilinéaire $\varphi$ sur un module $E$ de rang $3$ muni d'une base $\omega$ de
    $\det E$ définit un crochet sur le dual $E'$ ; $\varphi$ symétrique donne une
    algèbre de Lie. L'application cofacteur $f' = \Lambda^2 f$ donne un crochet
    sur $E$, $f : E \to E'$ est un homomorphisme, $E'$ opère sur $E$. L'exemple
    $xy + z^2$ donne $\mathfrak{sl}_2 \to \mathfrak{pgl}_2$. Formes de Casimir,
    changements d'échelle, caractéristique $2$. Puis $G = \mathrm{SO}(q)$, son
    revêtement $\widetilde{G}$, et l'isomorphisme $E' \simeq \mathrm{Lie}\, G$.
    La suite s'arrête sur une parenthèse qu'elle n'ouvre pas.
  \item \textbf{II. Le deuxième cahier (pages 33 à 55).} Une rotation d'un
    polygone régulier (page 33) ; la chaîne algèbres spéciales $=$ droites du plan
    $=$ diviseurs de degré $2$ de la conique $=$ sous-groupes commutatifs lisses
    de dimension $1$ (page 34) ; commutants et centralisateurs dans
    $\mathrm{PGL}(V)$, avec les contre-exemples de caractéristique $2$ (pages 36
    à 44) ; algèbres de quaternions, la conique $X \simeq \check{\mathbb{P}}(V)$,
    l'involution canonique, les sous-groupes « jacobiens » (pages 46 à 55).
  \item \textbf{III. Dictionnaire lié à $\mathrm{GP}(1) \simeq \mathrm{SO}(3)$
    (pages 57 à 78).} Neuf catégories fibrées, qui sont des gerbes ; leurs
    représentants typiques sur $\mathbf{Z}$ ; les opérations de
    $G_0 = \mathrm{PGL}_2$ ; les foncteurs du dictionnaire (III à V) ; enfin (VI)
    l'algèbre de quaternions comme clef de tous les objets. Le dossier s'arrête
    au milieu de VI, sur le morphisme de Frobenius.
\end{itemize}

\subsection*{Conventions, valables pour tout le dossier}

$S$ est un schéma ; « Module » veut dire $\mathcal{O}_S$-module localement libre
de rang fini, et un fibré vectoriel est identifié à son Module de sections.
$E$ et $E'$ sont de rang $3$, en dualité par $\langle\ ,\ \rangle$ ; $\omega$ est
une base de $\det E = \Lambda^3 E$, $\omega'$ la base duale de $\det E'$.
Une forme bilinéaire $\varphi$ sur $E$ et l'application $f : E \to E'$ se
déterminent l'une l'autre par $\langle y, f(x) \rangle = \varphi(x, y)$. Le
crochet construit sur $E'$ est noté $[\ ,\ ]_{\varphi}$, celui construit sur $E$
est noté $[\ ,\ ]_{\varphi'}$, comme sur les pages. Une forme quadratique $q$ et
sa forme polaire $\varphi$ sont liées par $q(x + y) = q(x) + q(y) + \varphi(x, y)$,
d'où $\varphi(x, x) = 2q(x)$.

Pour un Module $F$, $\check{\mathbb{P}}(F)$ est le schéma des sous-fibrés de rang $1$
de $F$ et $\mathbb{P}(F)$ celui de ses hyperplans (sous-fibrés de corang $1$),
de sorte que $\mathbb{P}(F) = \check{\mathbb{P}}(F^{\vee})$ ; c'est la convention des
pages, que celles-ci ne tiennent pas toujours (page 34). En rang $2$ les deux
coïncident.

$\mathrm{GP}(1)$ est le groupe projectif $\mathrm{PGL}_2$, $\mathfrak{gp}(1)$ son
algèbre de Lie $\mathfrak{pgl}_2 = \mathfrak{gl}_2/\mathcal{O}\cdot 1$ ; on garde la
notation des pages quand on les cite. $G = \mathrm{SO}(q)$ est le groupe des
automorphismes de $(E, q, \omega)$, $\widetilde{G}$ son revêtement simplement connexe,
$\mathfrak{g}$ et $\widetilde{\mathfrak{g}}$ leurs algèbres de Lie. Pour une algèbre de
quaternions $\mathcal{M}$, $\mathcal{M}^{\times}$ est son groupe des unités (les pages
écrivent $\check{\mathbb{V}}(\mathcal{M})^{*}$ ou $\mathcal{L}(\mathcal{M})$), $\det$ et
$\mathrm{Tr}$ ses norme et trace réduites.

\textbf{Deux homonymies.} Les lettres $X, Y, H$ désignent, dans la première
suite et à la page 40, une base de $\mathfrak{gp}(1)$ ; à partir de la page 34, $X$
est une conique et $Y$ un diviseur de degré $2$ sur elle. Les deux usages ne se
croisent jamais dans une même section. D'autre part, les pages 36 à 43 nomment
$H$ un sous-groupe commutatif lisse de dimension relative $1$ ; c'est le même
objet que le sous-groupe « jacobien » $\mathcal{T}$ de la page 54, et on l'écrit
$\mathcal{T}$ partout ici. De même $\Delta$ est, dans la première suite, le scalaire
déterminant de $f$ ; la forme $u \mapsto -\det u$ sur une algèbre de quaternions,
que la page 47 nomme aussi $\Delta$, est écrite ici $Q_{\mathcal{M}}$, et la quadrique
qu'elle définit $\mathcal{Q}$.

\subsection*{Ce que seule une lecture d'ensemble peut dire}

\begin{itemize}
  \item La page 31 annonce un isomorphisme $E \simeq \widetilde{\mathfrak{g}}$ et dit
    qu'il faut pour cela « redéfinir avec précision, en termes élémentaires, la
    construction de $\widetilde{G}$ ». La section VI du Dictionnaire (pages 68 à
    75) le fait en substance : $\widetilde{G}$ est le noyau de la norme réduite d'une
    algèbre de quaternions, $\widetilde{\mathfrak{g}}$ ses éléments de trace nulle,
    $q = -\det$, et $(E, q, \omega) = (\widetilde{\mathfrak{g}}, q, \omega)$ par
    définition (50), le signe de $\omega$ étant fixé par (52).
  \item Le signe de la page 19 ($f'(X) = 2\widetilde{X}$, etc.) est contredit par le
    dossier lui-même : la page 75 trouve $ff' = -2$, ce que donne aussi la formule
    de Cramer avec $\Delta = -2$. La relation de Casimir (74), qui en dépendait, se
    corrige en conséquence, et sa seconde moitié est la relation avec la forme de
    Killing que la page 72 énonce indépendamment.
  \item La dimension $6$ du groupe des automorphismes d'une forme de
    $\mathfrak{sl}(2)$ en caractéristique $2$ (page 27) est la raison de la restriction
    $(*)$ des pages 57 à 59 : le foncteur « algèbre de Lie » n'est une équivalence
    qu'au-dessus de $\mathbf{Z}[\tfrac{1}{2}]$.
  \item Le sous-groupe $H$ des pages 36 et 37, image de
    $(\mathcal{O}_S + \mathcal{O}_S u)^{\times}$, est le sous-groupe jacobien $\mathcal{T}$
    attaché à l'algèbre spéciale $\mathcal{A} = \mathcal{O}_S + \mathcal{O}_S u$ à la page 54 ;
    et le contre-exemple de caractéristique $2$ de la page 40 limite
    l'affirmation de la page 55 selon laquelle $\mathcal{T}$ est le centralisateur de
    son algèbre de Lie.
  \item La chaîne d'isomorphismes de la page 34 est le Corollaire (21)--(22) des
    pages 52 et 53, écrit avant sa démonstration dans l'ordre des feuillets.
\end{itemize}

\subsection*{Ce que le dossier annonce sans l'établir}

L'isomorphisme $E \simeq \widetilde{\mathfrak{g}}$ pour un module spécial quadratique
quelconque, avec ses propriétés (page 31) ; les équivalences de la page 46 et
les isomorphismes (7) de la page 62, admis « pour ne pas nous éterniser » ; le
dictionnaire complet de la page 58 ; la démonstration de l'énoncé sur les
sous-groupes d'ordre premier (pages 41 à 43) ; la question de savoir si un
sous-groupe lisse de dimension $1$ sur une base est automatiquement commutatif,
et le normalisateur (pages 43 et 44, qui s'interrompent) ; la question en marge
de la page 77. Le dossier s'arrête à la page 78, au milieu de VI.

\pagerange{2}{31}
\section*{I. Algèbres de Lie de rang $3$ associées aux formes bilinéaires (pages 1 à 31)}

\pagerange{2}{3}
\subsection*{La donnée, et les deux isomorphismes $\alpha$, $\alpha'$ (pages 2 et 3)}

Soient $E$, $E'$ deux Modules de rang $3$ sur $S$ en dualité parfaite
$E \times E' \to \mathcal{O}_S$ ; d'où $\det E \otimes \det E' \simeq \mathcal{O}_S$. Les
accouplements $E \times \Lambda^2 E \to \det E$ et $E' \times \Lambda^2 E' \to \det E'$
sont parfaits, d'où des isomorphismes
$\Lambda^2 E \simeq E' \otimes \det E$ et $\Lambda^2 E' \simeq E \otimes \det E'$. On
suppose $E$ \emph{spécial}, c'est-à-dire muni d'une base $\omega$ de $\det E$ ; soit
$\omega'$ la base duale de $\det E'$. Les isomorphismes précédents deviennent
\[
  \alpha : \Lambda^2 E \xrightarrow{\ \sim\ } E', \qquad
  \alpha' : \Lambda^2 E' \xrightarrow{\ \sim\ } E,
\]
caractérisés par
\[
  (7) \qquad \langle x, \alpha(\psi) \rangle\, \omega = x \wedge \psi , \qquad
  \langle \alpha'(\psi'), x' \rangle\, \omega' = \psi' \wedge x'
\]
pour $x \in E$, $\psi \in \Lambda^2 E$, $x' \in E'$, $\psi' \in \Lambda^2 E'$\footnote{La
page écrit le couple $(x, \omega)$ et son image $x \wedge \varphi$ : la lettre
$\omega$, réservée ensuite à la base de $\det E$, est un lapsus, de même que
$\varphi$ pour un élément de $\Lambda^2 E$. On écrit $\psi$. La lecture
« $\omega' = \omega^{-1}$ » de (4') est douteuse sur la page ; c'est la base duale.}.
Si $(e_1, e_2, e_3)$ est une base de $E$ avec $\omega = e_1 \wedge e_2 \wedge e_3$, et
$(e'_i)$ la base duale, on pose $\bar{e}_1 = e_2 \wedge e_3$,
$\bar{e}_2 = e_3 \wedge e_1$, $\bar{e}_3 = e_1 \wedge e_2$, de même $\bar{e}'_i$, et
\[
  (10) \qquad \alpha(\bar{e}_i) = e'_i , \qquad \alpha'(\bar{e}'_i) = e_i .
\]
En coordonnées, $\alpha$ et $\alpha'$ ne sont autres que le produit vectoriel :
$\alpha(x \wedge y) = x \times y$ quand on identifie $E$ et $E'$ à
$\mathcal{O}_S^3$ par des bases duales.

\textbf{Observation.} On a des bijections canoniques entre les formes
bilinéaires $\varphi$ sur $E$, les applications linéaires $f : E \to E'$ (par
$\langle y, f(x) \rangle = \varphi(x, y)$) et les structures d'algèbre alternée
sur $E'$, la dernière associant à $f$ le crochet
\[
  (12) \qquad [x', y']_{\varphi} = f\bigl(\alpha'(x' \wedge y')\bigr) .
\]
La bijection est celle de $\mathrm{Hom}(E, E') \simeq \mathrm{Hom}(\Lambda^2 E', E')$
déduite de $\alpha'$ : toute algèbre alternée sur un module de rang $3$ s'écrit
ainsi, et c'est l'énoncé moderne qu'une algèbre de Lie de dimension $3$ est
codée par un tenseur de type $(1,1)$\footnote{Les pages dessinent ces bijections
comme un triangle passant par $\mathrm{Hom}_{\mathcal{O}_S}(E, E')$. La formule (11) qui
définit $f$ porte à la page 3 un numéro douteux, et le numéro (11) sert de
nouveau à la page 5.}.

\pagerange{5}{6}
\subsection*{L'identité de Jacobi (pages 5, 4 et 6)}

Les pages 4 et 5 se lisent à rebours : la page 5 précède la page 4, qui commence
au milieu de sa dernière phrase. Dans les bases précédentes, notons
$\varphi_{ij} = \varphi(e_i, e_j)$ ; la matrice de $f$ est la transposée de
$(\varphi_{ij})$, et (12) donne
\[
  (12) \qquad [e'_2, e'_3]_{\varphi} = \sum_j \varphi_{1j}\, e'_j , \quad
  [e'_3, e'_1]_{\varphi} = \sum_j \varphi_{2j}\, e'_j , \quad
  [e'_1, e'_2]_{\varphi} = \sum_j \varphi_{3j}\, e'_j .
\]
Les symétries apparaissent en écrivant
\[
  (15) \qquad (\varphi_{ij}) = \begin{pmatrix} a & r & q' \\ r' & b & p \\ q & p' & c \end{pmatrix},
\]
de sorte que la transposition échange $p \leftrightarrow p'$, $q \leftrightarrow q'$,
$r \leftrightarrow r'$ et qu'une permutation circulaire de la base permute
circulairement $(a, b, c)$, $(p, q, r)$, $(p', q', r')$. On trouve
\[
  (17) \qquad
  \sum_{\text{circ}} [[e'_1, e'_2], e'_3]
  = \bigl[a(p' - p) + q'r' - qr\bigr] e'_1
  + \bigl[b(q' - q) + r'p' - rp\bigr] e'_2
  + \bigl[c(r' - r) + p'q' - pq\bigr] e'_3 ,
\]
et $E'$ est une algèbre de Lie si et seulement si les trois coefficients sont
nuls — conditions (18) —, c'est-à-dire si trois différences de mineurs d'ordre
$2$ de $(\varphi_{ij})$ et de sa transposée s'annulent (19)\footnote{Calcul refait
pour cette lecture ; il donne exactement (16) et (17) de la page 4. La
disposition des indices des mineurs dans (19) est incertaine sur la page.}.

\textbf{Proposition (page 6).} \emph{Si $\varphi$ est symétrique, $[\ ,\ ]_{\varphi}$
est une structure d'algèbre de Lie sur $E'$.} En effet $p = p'$, $q = q'$,
$r = r'$ annulent (17)\footnote{La réciproque est fausse, et la page ne la
prétend pas (« on trouve en particulier ») : une forme alternée vérifie aussi
(18). Lorsque $2$ est inversible et qu'on écrit $\varphi = \sigma + \beta$, partie
symétrique et partie alternée, $\beta$ correspondant à un vecteur $v$, le
jacobien (17) vaut $-2\,\sigma(v)$ (calcul de cette lecture) ; la condition de
Jacobi est $\sigma(v) = 0$. C'est, sur $\mathbf{R}$, la classification de Bianchi
(1898) des algèbres de Lie de dimension $3$, dont les « unimodulaires » sont
exactement celles de $\varphi$ symétrique.}.

\pagerange{6}{6}
\subsection*{L'exemple : $xy + z^2$ et l'algèbre de Lie de $\mathrm{PGL}_2$ (page 6)}

Prenons $E = \mathcal{O}_S^3$, $\omega = e_1 \wedge e_2 \wedge e_3$ et
\[
  (21) \qquad q(x, y, z) = xy + z^2 , \qquad
  \varphi\bigl((x, y, z), (x', y', z')\bigr) = xy' + yx' + 2zz' ,
\]
donc $a = b = 0$, $c = 2$, $p = p' = q = q' = 0$, $r = r' = 1$\footnote{La page
écrit $c = 1$ dans (23 bis), alors que la matrice qu'elle donne à côté porte
$\varphi_{33} = 2$, ce qu'impose (22). On corrige.}, et $f(e_1) = e'_2$,
$f(e_2) = e'_1$, $f(e_3) = 2e'_3$. Les formules (12) donnent
$[e'_2, e'_3] = e'_2$, $[e'_3, e'_1] = e'_1$, $[e'_1, e'_2] = 2e'_3$, c'est-à-dire,
avec $X = e'_1$, $Y = e'_2$, $H = e'_3$,
\[
  (26) \qquad [H, X] = X , \qquad [H, Y] = -Y , \qquad [X, Y] = 2H ,
\]
les formules de l'algèbre de Lie de $\mathrm{GP}(1) = \mathrm{PGL}_2$ sur $\mathbf{Z}$,
« universelle » au sens où elle est définie sur $\mathbf{Z}$\footnote{Avec
$X = E_{12}$, $Y = E_{21}$ et $H$ la classe de $E_{11}$ dans
$\mathfrak{gl}_2/\mathcal{O}\cdot 1$, on retrouve (26) : $[X, Y] = E_{11} - E_{22}$, qui est
congru à $2E_{11}$ modulo les scalaires. Ce n'est pas l'algèbre
$\mathfrak{sl}_2$ : en caractéristique $2$, les deux diffèrent, et c'est tout
l'objet des pages 21 et 26.}.

\pagerange{7}{8}
\subsection*{Discriminant et structure spéciale orthogonale (pages 7 et 8)}

Le discriminant d'une forme bilinéaire $\varphi$ sur $E$ est
$\delta(\varphi) = \det(\varphi_{ij})\, \omega'^{\otimes 2}$, section de
$(\det E^{\vee})^{\otimes 2} = (\det E)^{\otimes -2}$, indépendante de la base\footnote{Un
trait précède « det » dans (27) ; ce n'est pas un signe moins, comme le montre
le calcul de (30).}. Pour la forme polaire d'une forme quadratique
$q = ax^2 + by^2 + cz^2 + pyz + qzx + rxy$, de matrice
$\begin{pmatrix} 2a & r & q \\ r & 2b & p \\ q & p & 2c \end{pmatrix}$,
\[
  (30) \qquad \delta(\varphi) = 2\,(4abc + pqr - ap^2 - bq^2 - cr^2)\, \omega'^{\otimes 2} ,
\]
et l'on pose $\delta'(q) = (4abc + pqr - ap^2 - bq^2 - cr^2)\, \omega'^{\otimes 2}$, le
\emph{discriminant divisé} (aujourd'hui : demi-discriminant), qui garde un sens
en caractéristique $2$. Pour $q = xy + z^2$, $\delta'(q) = -\omega'^{\otimes 2}$.

D'où la définition : une \emph{structure spéciale orthogonale} sur un Module $E$
de rang $3$ est un couple $(q, \omega)$, $q$ forme quadratique et $\omega$ base de
$\det E$, tel que
\[
  (33) \qquad \delta'(q) = -\,\omega^{\otimes(-2)} .
\]
En rang $2n$ ou $2n + 1$, on demanderait $\delta'(q) = (-1)^n \omega^{\otimes(-2)}$
(où $\delta'$ est le discriminant en rang pair), normalisation choisie pour que la
forme standard $\sum_{i=1}^{n} x_{2i-1}x_{2i}$, resp. cette forme plus
$x_{2n+1}^2$, l'admette avec $\omega = e_1 \wedge \cdots \wedge e_N$\footnote{La page
indexe la somme par $x_{2i}x_{2i+1}$, $i = 1, \ldots, n$, ce qui décale les indices ;
le sens est clair.}. Une structure spéciale orthogonale associée à $q$ existe
localement si et seulement si $q$ est lisse, c'est-à-dire $\delta'(q)$ inversible,
et ces structures forment un torseur sous $\mu_2$ ($\omega \mapsto -\omega$)\footnote{
« Localement » doit s'entendre pour la topologie fppf : il faut extraire une
racine carrée de l'unité $-\delta'(q)$, ce qui n'est possible étale-localement
que si $2$ est inversible. C'est aussi pour la topologie fppf que $\mu_2$ est
un faisceau en groupes non trivial en caractéristique $2$. La fin de la phrase
sur la page est en partie illisible.}.

La construction précédente s'applique en échangeant $E$ et $E'$ : une forme
bilinéaire sur $E'$, c'est-à-dire une application $E' \to E$, définit une algèbre
alternée sur $E$.

\pagerange{9}{13}
\subsection*{L'application cofacteur $f'$ et le crochet sur $E$ (pages 9 à 13)}

À $f : E \to E'$ on associe $\Lambda^2 f$, d'où
\[
  (38) \qquad f' = \alpha' \circ \Lambda^2 f \circ \alpha^{-1} : E' \longrightarrow E ,
  \qquad f'\bigl(\alpha(x \wedge y)\bigr) = \alpha'\bigl(f(x) \wedge f(y)\bigr)
\]
\footnote{La page écrit $\alpha \circ \Lambda^2 f \circ \alpha^{-1}$ ; (6) et (39) imposent
$\alpha'$ à gauche.}. Dans des bases duales, la matrice de $f'$ est formée des
neuf mineurs d'ordre $2$ de celle de $f$ : c'est la \emph{comatrice}, et $f'$ est
symétrique avec $f$. La forme $\varphi'$ sur $E'$ associée à $f'$ définit donc,
par la construction échangée, une structure d'\emph{algèbre de Lie} sur $E$ :
\[
  (40\ \text{bis}) \qquad [x, y]_{\varphi'} = f'\alpha(x \wedge y) = \alpha'\bigl(f(x) \wedge f(y)\bigr),
  \qquad \bigl\langle [x, y]_{\varphi'}, x' \bigr\rangle\, \omega' = f(x) \wedge f(y) \wedge x' .
\]

\textbf{Cramer.} Écrivons $\delta(\varphi) = \Delta\, \omega^{\otimes -2}$, $\Delta \in
\Gamma(S, \mathcal{O}_S)$ ; $\Delta$ est le déterminant de $f$ relativement à
$\omega$, $\omega'$. Les formules de Cramer donnent
\[
  (42) \qquad f'f = \Delta\, \mathrm{id}_E , \qquad ff' = \Delta\, \mathrm{id}_{E'} .
\]
La relation entre $f$ et $f'$ n'est pas symétrique : $f' = 0$ si et seulement si
$f$ est de rang $\leq 1$ en chaque point, alors que $f$ de rang $2$ donne $f'$ de
rang $1$\footnote{La fin de la phrase de la page 10 est illisible ; on en donne
le contenu mathématique, qui est l'asymétrie annoncée.}. Plus précisément, avec
$\delta(\varphi') = \Delta' \omega'^{\otimes 2}$, on a $f'f'' = \Delta'\,\mathrm{id}$,
$f''f' = \Delta'\,\mathrm{id}$, et lorsque $\Delta$ est inversible
\[
  (45) \qquad f' = \Delta f^{-1} , \qquad \det f' = \Delta^2 , \qquad \Delta' = \Delta^2 ,
\]
d'où
\[
  (48) \qquad f'' = \Delta\, f ,
\]
identité qui, par le « principe de conservation des identités algébriques »,
vaut pour tout $f$ : c'est l'identité classique $\mathrm{adj}(\mathrm{adj}\, M) =
(\det M)^{n-2} M$ en dimension $n = 3$\footnote{Les numéros de la page 12 sont
surchargés ((46) ou (47)) ; on garde (45) et (48).}. Si $\Delta = 1$, $f$ et $f'$
sont inverses l'un de l'autre et la relation redevient symétrique. Mais ce qui
intéressera, dit la page 13, c'est l'exemple (21), où $\Delta = -2$, sans
s'interdire de spécialiser en caractéristique $2$, où $\Delta = 0$\footnote{La
page 13 écrit $\Delta = 2$ ; la matrice (23 bis) a pour déterminant $-2$, valeur
que la page 18 retrouve.}.

\textbf{Changements d'échelle (pages 10 et 11).} Si $\omega$ est remplacé par
$\lambda\omega$, $\lambda$ inversible, $\varphi$ fixé, alors $\alpha$ est divisé et
$\alpha'$ multiplié par $\lambda$ ; $\Delta$ et $f'$ sont multipliés par $\lambda^2$,
le crochet $[\ ,\ ]_{\varphi}$ sur $E'$ par $\lambda$ ; le crochet sur $E$ défini par
une $\varphi'$ \emph{fixée} serait multiplié par $\lambda^{-1}$, mais comme $\varphi'$
est elle-même multipliée par $\lambda^2$, le crochet $[\ ,\ ]_{\varphi'}$ sur $E$ est
multiplié par $\lambda$. Les deux crochets attachés à $(\varphi, \omega)$ varient donc
du même facteur $\lambda$\footnote{La page attribue $[\ ,\ ]_{\varphi}$ à $E$ et
$[\ ,\ ]_{\varphi'}$ à $E'$, au rebours de (12) et (40 bis) ; on rétablit.}. À
$\omega$ fixé, $\varphi \mapsto \mu\varphi$ multiplie $[\ ,\ ]_{\varphi}$ par $\mu$ et
$[\ ,\ ]_{\varphi'}$ par $\mu^2$ : le premier crochet dépend linéairement de $\varphi$,
le second quadratiquement.

\pagerange{13}{14}
\subsection*{$f$ est un homomorphisme ; $f'$ ne l'est pas (pages 13 et 14)}

\textbf{Programme (page 13).} On va montrer que $f : E \to E'$ est un
homomorphisme d'algèbres, faire opérer (si $\varphi$ est symétrique) l'algèbre de
Lie $E'$ sur $E$ par dérivations qui préservent $\varphi$, de façon que $f$ commute
aux actions de $E'$ sur $E$ et sur $E'$, et, lorsque $\varphi$ provient d'une forme
quadratique lisse spéciale, interpréter le tout par un homomorphisme de schémas
en groupes lisses $\widetilde{G} \to G$ et une opération de $G$ sur $\widetilde{G}$.

\textbf{$f$ est un homomorphisme.} Pour tout $\varphi$, symétrique ou non,
\[
  (49) \qquad f[x, y]_{\varphi'} = [f(x), f(y)]_{\varphi} :
\]
le second membre vaut $f\alpha'(f(x) \wedge f(y))$ par (12), et
$[x, y]_{\varphi'} = \alpha'(f(x) \wedge f(y))$ est (40 bis).

\textbf{$f'$ n'en est pas un, sauf si $\Delta = 1$.} Supposons $\Delta$ inversible.
Alors $f^{-1}$ est un isomorphisme d'algèbres, et $f' = \Delta f^{-1}$ ; si $f'$
était un homomorphisme, on aurait à la fois
$f^{-1}[x', y'] = [f^{-1}x', f^{-1}y']$ et
$\Delta f^{-1}[x', y'] = \Delta^2 [f^{-1}x', f^{-1}y']$, donc
$(\Delta^2 - \Delta)[f^{-1}x', f^{-1}y'] = 0$ ; comme les crochets engendrent $E$
(le crochet $[\ ,\ ]_{\varphi'}$ est surjectif, $f'$ et $\alpha$ l'étant), $\Delta^2 = \Delta$,
soit $\Delta = 1$\footnote{La page conclut « dès que $\Delta \neq 1$, $\Delta$
inversible, $f'$ n'est pas un hom. d'algèbres » ; sur une base non connexe, il
faut lire : $f'$ est un homomorphisme si et seulement si $\Delta = 1$. Une
lecture de la page 14 donne $\Delta' b$ au lieu de $\Delta^2 b$.}.

\pagerange{15}{18}
\subsection*{L'action de $E'$ sur $E$ (pages 15, 17, 16 et 18)}

L'argument court de la page 15 à la page 17, puis 16, puis 18 : chaque page
commence au milieu de la phrase qui achève la précédente dans cet ordre. On
suppose désormais $\varphi$ symétrique.

\textbf{L'énoncé (page 15).} À tout $\xi' \in E'$ on veut associer un endomorphisme
$\rho_E(\xi')$ de $E$ qui préserve $\varphi$ infinitésimalement,
\[
  (52) \qquad \varphi\bigl(\rho_E(\xi')x, y\bigr) + \varphi\bigl(x, \rho_E(\xi')y\bigr) = 0 ,
\]
qui soit une dérivation du crochet $[\ ,\ ]_{\varphi'}$ de $E$ (53), tel que
$\rho_E : E' \to \underline{\mathrm{End}}(E)$ soit un homomorphisme d'algèbres de Lie,
\[
  (55) \qquad \rho_E\bigl([\xi', \eta']_{\varphi}\bigr) = [\rho_E(\xi'), \rho_E(\eta')] ,
\]
et que $f$ entrelace $\rho_E$ et la représentation adjointe de $E'$ :
\[
  (56) \qquad f\bigl(\rho_E(\xi')x\bigr) = [\xi', f(x)]_{\varphi}
\]
\footnote{Dans (55), (57) et (59 bis), l'indice du crochet de $E'$ est écrit
$\varphi'$, ou d'un signe épais qui peut être un astérisque ; le crochet de $E'$
est $[\ ,\ ]_{\varphi}$ (12), comme le confirme l'accolade de (59 bis),
$f\alpha'(\xi' \wedge x')$.}. Que (52) entraîne (53) est un « transport de
structure » : $1 + \varepsilon\rho_E(\xi')$ est un automorphisme de $E$ sur les nombres
duaux $S[\varepsilon]$, et le crochet $[\ ,\ ]_{\varphi'}$ est fonctoriel en
$(\varphi, \omega)$. Il faut pour cela que cet automorphisme préserve aussi $\omega$,
c'est-à-dire
\[
  (60) \qquad \mathrm{Tr}\, \rho_E(\xi') = 0 ,
\]
que la page ne note qu'en marge de la page 18 ; c'est vrai ici, parce qu'une
algèbre de Lie de la forme (12) avec $\varphi$ symétrique est unimodulaire
($\mathrm{Tr}\, \mathrm{ad}\, \xi' = 0$)\footnote{Hypothèse fournie par l'édition :
la page 15 dit que (52) « pour toute base » implique (53). En caractéristique
$\neq 2$, la nullité de la trace découle de (52) pour $\varphi$ non dégénérée ; en
général elle découle du calcul.}.

\textbf{La construction (page 17).} $E'$ opère sur lui-même par la représentation
adjointe $\rho_{E'}(\xi')x' = [\xi', x']_{\varphi}$ ; on la fait opérer sur
$E = (E')^{\vee}$ par la contragrédiente,
\[
  (58) \qquad \rho_E(\xi') = -\,{}^{t}\rho_{E'}(\xi'),
\]
c'est-à-dire
\[
  (59\ \text{bis}) \qquad \bigl\langle \rho_E(\xi')x, x' \bigr\rangle
  = -\bigl\langle x, f\alpha'(\xi' \wedge x') \bigr\rangle
  = -\varphi\bigl(\alpha'(\xi' \wedge x'), x\bigr).
\]
La contragrédiente d'une représentation est une représentation : c'est (55).

\textbf{La formule précise (page 16).} On a
\[
  (56\ \text{bis}) \qquad \rho_E(\xi')x = \alpha'\bigl(\xi' \wedge f(x)\bigr).
\]
En effet, multipliant (59 bis) par $\omega'$ et utilisant (7) et la symétrie de
$\varphi$, il s'agit de voir $-\,\xi' \wedge x' \wedge f(x) = \xi' \wedge f(x) \wedge x'$,
ce qui est vrai — « hurrah ! », écrit la page. Appliquant $f$, on obtient (56).
Enfin (52) s'écrit
$\langle y, f(\rho_E(\xi')x) \rangle + \langle \rho_E(\xi')y, f(x) \rangle = 0$, soit,
grâce à (56), $\langle y, \rho_{E'}(\xi')f(x) \rangle + \langle \rho_E(\xi')y, f(x) \rangle = 0$,
ce qui est la définition (58)--(59) appliquée à $x' = f(x)$ (page 18).

\textbf{La formule « oubliée » (marge de la page 15).} Soit $\tilde{\rho}_E$ la
représentation adjointe de $E$ sur lui-même,
$\tilde{\rho}_E(x)y = [x, y]_{\varphi'}$. Alors
\[
  \rho_E\bigl(f(x)\bigr) = \tilde{\rho}_E(x) ,
\]
c'est-à-dire que l'action de $E'$, restreinte à $E$ à travers $f$, est l'action
adjointe : par (56 bis) et (40 bis),
$\rho_E(f(x))y = \alpha'(f(x) \wedge f(y)) = [x, y]_{\varphi'}$\footnote{La note
marginale porte aussi le numéro (56 bis) ; elle est écrite en remontant le long
du bord, et la place des accents dans « $\rho_E(f(x')) = \tilde{\rho}_E(x')$ » est
incertaine. Ce qu'on énonce est ce que (56 bis) de la page 16 entraîne.}.

\textbf{La forme quadratique (marge de la page 18).} Faisant $x = y$ dans (52), on
trouve $2\varphi(x, \rho_E(\xi')x) = 0$ ; en fait $\varphi(x, \rho_E(\xi')x) = 0$
(52 bis), par prolongement à partir du cas universel sur $\mathbf{Z}$, où $2$ n'est
pas diviseur de zéro. Cela signifie que $\rho_E(\xi')$ préserve infinitésimalement
la \emph{forme quadratique} $q$, et non seulement sa forme polaire — distinction
qui n'a de contenu qu'en caractéristique $2$\footnote{La marge de la page 18 est
en grande partie illisible ; l'indice de $\rho$ y est écrit $E'$. On donne la
seule lecture qui ait un sens.}.

\pagerange{18}{21}
\subsection*{L'exemple explicité : $\mathfrak{sl}_2 \to \mathfrak{pgl}_2$ (pages 18 à 21)}

Revenons à (21) : $E'$ a la base $X, Y, H$ de (26), et
$\omega' = e'_1 \wedge e'_2 \wedge e'_3 = -X \wedge H \wedge Y$\footnote{La page note
qu'on aurait pu prendre l'autre signe « sans inconvénient » : $(q, -\omega)$ est
encore une structure spéciale orthogonale, les deux crochets changeant alors
de signe.}. Posons
\[
  (62) \qquad \widetilde{X} = e_2 , \quad \widetilde{Y} = e_1 , \quad \widetilde{H} = e_3 ,
  \qquad \text{d'où} \qquad \omega = \widetilde{X} \wedge \widetilde{H} \wedge \widetilde{Y}
\]
(« il y a un signe ici ! »), et
\[
  (65) \qquad f(\widetilde{X}) = X , \quad f(\widetilde{Y}) = Y , \quad f(\widetilde{H}) = 2H ,
  \qquad \Delta = -2 .
\]
Alors
\[
  (66) \qquad f'(X) = -2\widetilde{X} , \quad f'(Y) = -2\widetilde{Y} , \quad f'(H) = -\widetilde{H} ,
\]
conformément à $f'f = \Delta = -2$\footnote{La page écrit $f'(X) = 2\widetilde{X}$,
$f'(Y) = 2\widetilde{Y}$, $f'(H) = \widetilde{H}$, soit $f'f = +2$, en contradiction
avec (42) et $\Delta = -2$. Le calcul direct, $f'(H) = \alpha'(f(e_1) \wedge f(e_2)) =
\alpha'(e'_2 \wedge e'_1) = -e_3$, donne le signe moins, et la page 75 du dossier
le retrouve : $ff' = -2$ (54). Les formules (67) à (69) qui en découlent changent
de signe avec elle.}. La forme $\varphi'$ sur $E'$ est donc, pour
$\xi = xX + hH + yY$ et $\xi' = x'X + h'H + y'Y$,
\[
  (68) \qquad \varphi'(\xi, \xi') = -\bigl(hh' + 2(xy' + yx')\bigr) , \qquad
  \varphi'(\xi, \xi) = -(h^2 + 4xy) .
\]
Comme $\varphi'(H, H) = -1$ est impair, $\varphi'$ n'est la forme polaire d'aucune
forme quadratique quand $2$ n'est pas inversible : en caractéristique $2$,
$\varphi'(\xi, \xi) = h^2$ n'est pas nul, alors qu'une forme polaire y est alternée.
La page ajoute que, dans ces coordonnées, $q$, $\varphi$ et $\varphi'$ sont à
coefficients entiers premiers entre eux, donc non nulles sur tout corps.

Le crochet de $E$ est
\[
  (70) \qquad [\widetilde{H}, \widetilde{X}] = 2\widetilde{X} , \quad
  [\widetilde{H}, \widetilde{Y}] = -2\widetilde{Y} , \quad
  [\widetilde{X}, \widetilde{Y}] = \widetilde{H} ,
\]
celui de $\mathfrak{sl}(2)$, algèbre de Lie de $\mathrm{SL}(2)_S$, et $f$ est
l'homomorphisme d'algèbres de Lie de l'isogénie
\[
  (71) \qquad \mathrm{SL}(2)_S \longrightarrow \mathrm{GP}(1)_S
\]
de noyau $\mu_2$, centre de $\mathrm{SL}(2)$. L'action de $E'$ sur $E$ est la
différentielle de l'action de $\mathrm{GP}(1)$ sur $\mathrm{SL}(2)$ déduite par
passage au quotient de l'action adjointe, qui est triviale sur le centre. Avec
(56 bis), le plus commode :
\[
  (72) \qquad
  \begin{cases}
    \rho_E(X) = \tilde{\rho}_E(\widetilde{X}) : & \widetilde{X} \mapsto 0 ,\ \widetilde{Y} \mapsto \widetilde{H} ,\ \widetilde{H} \mapsto -2\widetilde{X} \\
    \rho_E(Y) = \tilde{\rho}_E(\widetilde{Y}) : & \widetilde{X} \mapsto -\widetilde{H} ,\ \widetilde{Y} \mapsto 0 ,\ \widetilde{H} \mapsto 2\widetilde{Y} \\
    \rho_E(H) : & \widetilde{X} \mapsto \widetilde{X} ,\ \widetilde{Y} \mapsto -\widetilde{Y} ,\ \widetilde{H} \mapsto 0
  \end{cases}
\]
où $\rho_E(H)$ est « $\tfrac{1}{2}\tilde{\rho}_E(\widetilde{H})$ (sic) » : la moitié d'un
endomorphisme à coefficients pairs, bien défini sur $\mathbf{Z}$.

\textbf{Remarque (page 21).} Si $2$ est inversible, $f$ est un isomorphisme
d'algèbres de Lie, qui identifie $E$ à $E'$ (l'identification usuelle d'un espace
quadratique avec son dual), $\rho_E$ à l'action adjointe, et $f' = \Delta f^{-1}$ à
$-2\,\mathrm{id}_E$, sans grand intérêt\footnote{La page écrit $2\,\mathrm{id}_E$,
suite du signe de (66).} ; de même chaque fois que $\Delta$ est inversible. Si $2$
est nul sur $S$, $f$ a pour noyau $\mathcal{O}_S\widetilde{H}$, qui est le centre de
$\mathfrak{sl}(2)$ et aussi son algèbre dérivée\footnote{La page écrit « groupe
dérivé » ; il s'agit de l'algèbre dérivée $[E, E]$. Le groupe dérivé de
$\mathrm{SL}_2$ est $\mathrm{SL}_2$ tout entier.} ; $f'$ a pour noyau
$\mathcal{O}_S X \oplus \mathcal{O}_S Y$, qui est l'algèbre dérivée de $\mathfrak{gp}(1)$ et
l'image de $f$, et l'image de $f'$ est le noyau $\mathcal{O}_S\widetilde{H}$ de $f$. Le
centre de $\mathfrak{gp}(1)$ est nul sur toute base. Il faudrait, dit la page,
expliciter les formes de $\mathfrak{sl}(2)$ et de $\mathfrak{gp}(1)$ en
caractéristique $2$ — ce que fait la page 26 — « contrairement aux formes de
$\mathrm{SL}(2)_S$ ou $\mathrm{GP}(1)_S$, qui sont équivalentes et équivalent aux formes
de $\mathbb{P}^1_S$ ».

\pagerange{22}{24}
\subsection*{Formes de Casimir (pages 22 à 24)}

Posons
\[
  (73) \qquad \mathrm{Cas}_{E'}(x', y') = \mathrm{Tr}\bigl(\rho_{E'}(x')\rho_{E'}(y')\bigr), \qquad
  \mathrm{Cas}_E(x, y) = \mathrm{Tr}\bigl(\tilde{\rho}_E(x)\tilde{\rho}_E(y)\bigr),
\]
les formes de Killing de $E'$ et de $E$\footnote{Un premier essai, définissant
des formes quadratiques $\mathrm{Tr}\,\rho(x)^2$, est barré en tête de la page 22.}.
Alors
\[
  (74) \qquad \mathrm{Cas}_{E'}(x', y') = -2\,\varphi'(x', y') , \qquad
  \mathrm{Cas}_E(x, y) = -2\Delta\, \varphi(x, y) .
\]
\footnote{La page écrit $\mathrm{Cas}_{E'} = 2\varphi'$, en cohérence avec son signe de
(66) ; avec le signe correct, c'est $-2\varphi'$. Les deux identités ont été
vérifiées pour une forme symétrique générale (calcul de cette lecture). La
seconde est juste sur la page.}
La page en donne une démonstration par réduction : les deux membres de chaque
identité ont le même comportement quand $\omega \mapsto \lambda\omega$ et
$\varphi \mapsto \mu\varphi$ (facteurs $\lambda^2\mu^2$ et $\lambda^2\mu^4$) ; par
conservation des identités polynomiales il suffit de les vérifier sur un corps
algébriquement clos de caractéristique $0$ pour une forme non dégénérée, que l'on
ramène, au prix de ces changements d'échelle, à (21) ; là, les deux membres sont
des formes symétriques invariantes sous le groupe spécial orthogonal, donc
proportionnelles, et il suffit de les comparer en un point :
$\mathrm{Cas}_{E'}(H, H) = 1 + 1 + 0 = 2$ et $\varphi'(H, H) = -1$ ;
$\mathrm{Cas}_E(\widetilde{H}, \widetilde{H}) = 4 + 4 + 0 = 8$ et
$-2\Delta\,\varphi(\widetilde{H}, \widetilde{H}) = -2 \cdot (-2) \cdot 2 = 8$\footnote{Les
quatre premières lignes de la page 24, qui font la réduction au cas (21), se
lisent mal ; seules les formules sont sûres.}.

Pour une structure spéciale orthogonale $(q, \omega)$, $\Delta = -2$ et
$\varphi(x, x) = 2q(x)$, d'où
\[
  (75)\text{--}(76) \qquad \mathrm{Cas}_E(x, x) = -4\Delta\, q(x) = 8\, q(x) .
\]
Donc, si $2$ est inversible (ou seulement non diviseur de zéro), $q$ \emph{se
déduit} de la structure d'algèbre de Lie de $E$ : c'est la forme de Killing
divisée par $8$. De même, la première formule redonne $\varphi'$ à partir du
crochet de $E'$, $\mathrm{Cas}_{E'}(x', x') = -2\varphi'(x', x')$ (77). C'est la
relation que la page 72 énonce, dans le langage des quaternions, entre $q = -\det$
et la forme de Killing de $\widetilde{\mathfrak{g}}$.

\pagerange{25}{25}
\subsection*{Quand deux données donnent les mêmes crochets (page 25)}

La page se présente comme un complément à la discussion de (42). On note
$[\ ,\ ]_{\omega, \varphi}$ le crochet (12) sur $E'$ et $[\ ,\ ]_{\omega, \varphi'}$ le
crochet (40 bis) sur $E$. D'après ce qui précède,
\[
  (78) \qquad [\ ,\ ]_{\lambda\omega, \mu\varphi} = \lambda\mu\, [\ ,\ ]_{\omega, \varphi} ,
  \qquad [\ ,\ ]_{\lambda\omega, (\mu\varphi)'} = \lambda\mu^2\, [\ ,\ ]_{\omega, \varphi'} .
\]
Comme, à $\omega$ fixé, $\varphi \mapsto [\ ,\ ]_{\varphi}$ est bijectif,
\[
  (79) \qquad [\ ,\ ]_{\omega, \varphi} = [\ ,\ ]_{\omega_1, \varphi_1}
  \iff \omega_1 = \lambda\omega ,\ \varphi_1 = \lambda^{-1}\varphi \ \text{pour un}\ \lambda \in \Gamma(S, \mathcal{O}_S^{\times}).
\]
Pour le crochet de $E$, on suppose $\Delta$ inversible. Alors, à $\omega$ fixé,
\[
  (80) \qquad f' = g' \iff g = \varepsilon f , \quad \varepsilon \in \Gamma(S, \mu_2) :
\]
si $f' = g'$, $\det f' = \det g'$ donne $\Delta(f)^2 = \Delta(g)^2$, donc
$\Delta(g) = \varepsilon\Delta(f)$ avec $\varepsilon^2 = 1$, et $f = \Delta(f)^{-1}f''$,
$g = \Delta(g)^{-1}g''$ donnent $g = \varepsilon f$\footnote{La page écrit d'abord
« $\Delta(f) = \Delta(g)$ i.e. $\Delta(f)^2 = \Delta(g)^2$ » ; c'est l'égalité des
déterminants de $f'$ et $g'$, $\det f' = \Delta^2$ par (45), qui donne celle des
carrés. L'hypothèse $\Delta$ inversible, implicite, est celle sous laquelle la
page vient de poser $f$ isomorphisme.}. On en déduit, toujours pour $\Delta$
inversible,
\[
  (81) \qquad [\ ,\ ]_{\omega_1, \varphi_1'} = [\ ,\ ]_{\omega, \varphi'}
  \iff \omega_1 = \mu^{-2}\omega ,\ \varphi_1 = \mu\varphi \ \text{pour un}\ \mu \in \Gamma(S, \mathcal{O}_S^{\times}),
\]
$\mu$ étant alors uniquement déterminé par $\varphi_1$\footnote{La page énonce :
$\omega_1 = \lambda\omega$, $\varphi_1 = \mu\varphi$ avec $\lambda\mu^2 = 1$, « $\mu$ une
racine carrée de $\lambda^{-1}$ ». C'est la même chose, à condition de lire que
$\lambda$ doit être l'inverse d'un carré : si $\omega_1 = \lambda\omega$ avec $\lambda$
non carré, aucun $\varphi_1$ ne donne le même crochet.}.

\pagerange{26}{27}
\subsection*{Caractéristique $2$ : algèbres couplées et plans affines symplectiques (pages 26 et 27)}

Supposons $2 = 0$ sur $S$. D'après la page 21, $\mathfrak{sl}(2)$ a un centre
$\mathfrak{z} = \mathcal{O}\widetilde{H}$ égal à son algèbre dérivée, le quotient
$V = \mathfrak{sl}(2)/\mathfrak{z}$ est de rang $2$ et le crochet induit un isomorphisme
$\Lambda^2 V \simeq \mathfrak{z}$ : c'est l'\emph{algèbre de Heisenberg} de dimension
$3$. Une forme de $\mathfrak{sl}(2)$ est donc la même chose qu'un couple $(V, E)$,
$V$ de rang $2$ et $E$ une extension de $V$ par $\det V$, le crochet étant
$[x, y] = \bar{x} \wedge \bar{y}$ ; pour $V$ fixé, une telle extension revient à une
extension de $\mathcal{O}_S$ par $(\det V) \otimes V^{\vee} \simeq V$, c'est-à-dire à un
fibré affine sous $V$. De même $\mathfrak{gp}(1)$ a l'idéal abélien
$V' = \mathcal{O}X \oplus \mathcal{O}Y$, de quotient $\mathcal{O}_S$, et tout relèvement de $1$
opère sur $V'$ par l'identité ; une forme de $\mathfrak{gp}(1)$ est un couple
$(V', E')$, $E'$ extension de $\mathcal{O}_S$ par $V'$\footnote{Sur la page, au-dessus
de $V'$, un $\check{V}'$ est écrit en interligne, qui peut le corriger ; la
lecture « Module vectoriel » est celle de la page 36.}.

Disons que deux telles algèbres $E$, $E'$ sont \emph{couplées} s'il y a entre elles
une dualité parfaite pour laquelle $\det V \subset E$ est l'orthogonal de
$V' \subset E'$. Alors $V' \simeq V^{\vee}$ et $\det V \simeq \mathcal{O}_S$ : $V$ reçoit une
structure symplectique et $V' \simeq V$. Les couples couplés équivalent donc aux
triples $(V, \varphi, E)$, $\varphi$ forme symplectique sur $V$ et $E$ fibré affine
sous $V$ : un fibré en \emph{plans affines symplectiques}. Le schéma en groupes de
leurs automorphismes est le groupe affine symplectique
$\mathrm{SL}(V) \ltimes V$, lisse de dimension relative $5$, tandis que le groupe
des automorphismes d'une forme de $\mathfrak{sl}(2)$, ou de $\mathfrak{gp}(1)$, est
une forme du groupe affine $\mathrm{GL}(V) \ltimes V$, lisse de dimension relative
$6$\footnote{La fin du premier paragraphe de la page 27 se lit mal, et ce qu'elle
compare à une forme de $\mathrm{GP}(1)$ reste incertain. La conséquence qu'on en
tire est de l'édition : comme $\mathrm{GP}(1)$ est de dimension $3$, le foncteur
« algèbre de Lie » ne peut être une équivalence au-dessus d'une base de
caractéristique $2$ ; c'est la restriction $(*)$ des pages 57 à 59.}.

Les applications $f$ et $f'$ s'en déduisent : $f$ est le composé
$E \to V \xrightarrow{\varphi} V' \to E'$, et $f'$ le composé
$E' \to \mathcal{O}_S \simeq \det V \to E$. Une description de $\rho_{E'}$ est commencée
puis barrée.

\pagerange{28}{31}
\subsection*{Le groupe $G = \mathrm{SO}(q)$, son revêtement, et $E' \simeq \mathfrak{g}$ (pages 28 à 31)}

Soient $q$ une forme quadratique lisse sur $E$, $\varphi$ sa forme polaire, et
$\omega$ une base de $\det E$ avec $\delta'(q) = -\omega^{\otimes(-2)}$ (82). Posons
\[
  (83) \qquad G = \mathrm{SO}(q) = \underline{\mathrm{Aut}}(E, q, \omega) ,
\]
et soit $\widetilde{G} \xrightarrow{F} G$ son revêtement simplement connexe, de noyau
canoniquement isomorphe à $\mu_2$ :
\[
  (84) \qquad 1 \longrightarrow \mu_2 \longrightarrow \widetilde{G} \xrightarrow{\ F\ } G \longrightarrow 1 .
\]
Aujourd'hui, $\widetilde{G}$ est le groupe $\mathrm{Spin}(q)$, le groupe des éléments
de norme $1$ de l'algèbre de Clifford paire de $q$, qui est une algèbre de
quaternions\footnote{Noms modernes, donnés par l'édition. Les pages écrivent
« revêtement simplement connexe » entre guillemets et ne nomment ni l'algèbre de
Clifford ni Spin ; la section VI (pages 68 et suivantes) construit $\widetilde{G}$
directement à partir d'une algèbre de quaternions. Pour $\mathrm{SO}(q)$ en rang
impair sur une base quelconque, $\underline{\mathrm{Aut}}(E, q, \omega)$ est la bonne
définition : $\mathrm{O}(q) = \mathrm{SO}(q) \times \mu_2$.}. Le noyau étant central,
l'action adjointe de $\widetilde{G}$ sur lui-même passe au quotient en une action
\[
  (87) \qquad \rho : G \longrightarrow \underline{\mathrm{Aut}}_{S\text{-}\mathrm{gr}}(\widetilde{G}) ,
\]
caractérisée par $\rho(F(a))(b) = aba^{-1}$ et $F(\rho(g)(b)) = g\,F(b)\,g^{-1}$
(diagrammes (88) et (89)). Sur les algèbres de Lie on obtient un homomorphisme
$\mathfrak{f} : \widetilde{\mathfrak{g}} \to \mathfrak{g}$ (90) et une action
$\rho_{\mathfrak{g}} : \mathfrak{g} \to \underline{\mathrm{Der}}(\widetilde{\mathfrak{g}})$ (91),
avec $\rho_{\mathfrak{g}}(\mathfrak{f}(a)) = \mathrm{ad}(a)$ et
$\mathfrak{f}(\rho_{\mathfrak{g}}(x)b) = [x, \mathfrak{f}(b)]$ (92).

\textbf{L'énoncé.} Il existe des isomorphismes canoniques
\[
  (93) \qquad i : E \xrightarrow{\ \sim\ } \widetilde{\mathfrak{g}} , \qquad
  j : E' \xrightarrow{\ \sim\ } \mathfrak{g}
\]
qui transportent les crochets de $\widetilde{\mathfrak{g}}$ et $\mathfrak{g}$ en
$[\ ,\ ]_{\varphi'}$ sur $E$ et $[\ ,\ ]_{\varphi}$ sur $E'$, l'homomorphisme
$\mathfrak{f}$ en $f$, et $\rho_{\mathfrak{g}}$ en $\rho_E$ ; et le composé
$\mathfrak{g} \to \underline{\mathrm{Der}}(\widetilde{\mathfrak{g}}) \simeq
\underline{\mathrm{Der}}(E) \subset \underline{\mathrm{End}}(E)$ est la
représentation de $\mathfrak{g}$ déduite de la représentation linéaire de $G$ sur
$E$.

\textbf{La moitié démontrée : $j$.} $\rho_E$ envoie $E'$ dans les endomorphismes de
$E$ qui préservent infinitésimalement $q$ (marge de la page 18) et $\omega$ (60),
lesquels forment exactement $\mathfrak{g} = \mathrm{Lie}\, \mathrm{SO}(q)$ ; d'où
$j : E' \to \mathfrak{g}$, compatible aux crochets par (55). C'est un isomorphisme :
$E'$ et $\mathfrak{g}$ étant localement libres de rang $3$, il suffit que $j$ soit
injectif sur chaque fibre, ce qui se voit après extension à un corps $k$
algébriquement clos. Là, on trouve une base où $q = xy + z^2$ et
$\omega = e_1 \wedge e_2 \wedge e_3$ (quitte à changer $e_3$ en $-e_3$), et (72) donne
\[
  \rho_E(xX + hH + yY) :
  \widetilde{X} \mapsto h\widetilde{X} - y\widetilde{H} , \quad
  \widetilde{Y} \mapsto x\widetilde{H} - h\widetilde{Y} , \quad
  \widetilde{H} \mapsto -2x\widetilde{X} + 2y\widetilde{Y} ,
\]
qui n'est nul que pour $x = y = h = 0$, en toute caractéristique\footnote{Le signe
devant $h\widetilde{Y}$ est surchargé sur la page 31 ; (72) donne $-$.}.

\textbf{La moitié annoncée : $i$.} « Il faut encore définir l'isom.
$E \simeq \widetilde{\mathfrak{g}}$ et prouver ses propriétés. Pour ceci il faut
redéfinir avec précision, en termes élémentaires, la construction de $\widetilde{G}$
sur $G$ — il nous faut pour ceci ouvrir une parenthèse. » La page s'arrête là, au
premier tiers de la feuille. La parenthèse n'est pas ouverte dans cette suite ;
elle l'est, en substance, par la section VI du Dictionnaire (pages 68 à 75), où
$\widetilde{G}$ est le noyau de la norme réduite d'une algèbre de quaternions et
où l'on pose $E = \widetilde{\mathfrak{g}}$.

\pagerange{33}{55}
\section*{II. Le deuxième cahier (pages 32 à 55)}

La couverture (page 32) ne porte que « (21 p) ». Le cahier n'est pas une suite
continue : une page de calcul sur les polygones, une page de résumé, une étude
des commutants dans $\mathrm{PGL}(V)$ en rang $2$, enfin les algèbres de quaternions.

\pagerange{33}{33}
\subsection*{La rotation d'un polygone régulier (page 33)}

En tête, à l'encre noire :
\[
  F_3 = T , \quad F_4 = T - 1 , \quad F_5 = T^2 - T - 1 , \quad F_6 = T - 2 ,
\]
$F_n$ étant le polynôme minimal sur $\mathbf{Q}$ de $1 + 2\cos(2\pi/n) = 1 + \zeta +
\zeta^{-1}$, $\zeta$ racine primitive $n$-ième de l'unité. Puis un polygone régulier
à $n$ sommets $\zeta_0, \zeta_1, \ldots$ inscrit dans un cercle, ses côtés
$e_1 = \zeta_1 - \zeta_0$, $e_2 = \zeta_2 - \zeta_1$, et la rotation affine $u$ qui
envoie $\zeta_i$ sur $\zeta_{i+1}$ :
\[
  u(\zeta_0) = \zeta_0 + e_1 , \qquad u(e_1) = e_2 , \qquad u(e_2) = -e_1 + (\alpha - 1)e_2 .
\]
La partie linéaire de $u$ a pour polynôme caractéristique $T^2 - 2\cos(2\pi/n)\,T +
1$, donc $\alpha - 1 = 2\cos(2\pi/n)$ ; et la flèche marquée $\alpha$ que la figure
trace de $\zeta_0$ à $\zeta_3$ s'explique : $\zeta_3 - \zeta_0 = e_1 + e_2 + u(e_2) =
\alpha\, e_2$. Ainsi $\alpha = 1 + 2\cos(2\pi/n)$ est le rapport de la diagonale
$\zeta_0\zeta_3$ au côté, racine de $F_n$ — le nombre d'or pour le pentagone, $2$
pour l'hexagone — et c'est aussi la trace de la rotation d'angle $2\pi/n$ de
l'espace à trois dimensions\footnote{Ces deux interprétations sont de l'édition ;
la page ne donne que les polynômes, la figure et les trois formules. Le
dossier 85 (pages 5 à 7, puis 23 à 31) travaille avec l'invariant
$2\cos\theta = \xi + \xi^{-1}$ d'une rotation et des polynômes
« ½cyclotomiques » $\Psi_n$ ; si $\Psi_n$ est, comme ce nom et les cas cités
($\alpha = -1, 0, 1$ et $\alpha^2 + \alpha - 1 = 0$ pour les ordres $3, 4, 6, 5$) le
suggèrent, le polynôme minimal de $2\cos(2\pi/n)$, alors le $F_n$ d'ici est
$\Psi_n(T - 1)$. Le même dossier 85 nomme aussi $F_n$ d'autres polynômes ;
l'homonymie est fortuite.}.

\pagerange{34}{34}
\subsection*{Algèbres spéciales, droites, diviseurs, sous-groupes (page 34)}

La page aligne :
\[
  \begin{aligned}
  \text{AlgSpéc}(\mathcal{M}) &\simeq \mathbb{P}(\widetilde{\mathfrak{g}}) \simeq
  \check{\mathbb{P}}(\mathfrak{g}) \simeq \check{P} = \mathrm{Dr}(P) \simeq \mathrm{Div}_2^{+}(X) \\
  &\simeq \{\text{sous-groupes commutatifs lisses de dimension } 1 \text{ de } G\} ,
  \end{aligned}
\]
c'est-à-dire : les sous-algèbres spéciales d'une algèbre de quaternions, les
droites du plan $P$, les diviseurs effectifs de degré $2$ de la conique $X \subset P$,
et les sous-groupes commutatifs lisses de dimension relative $1$ de $G$. C'est
l'énoncé que les pages 50 à 55 démontrent, auxquelles on renvoie\footnote{La page
écrit $\mathbb{P}(\mathfrak{g}) \simeq \check{\mathbb{P}}(\widetilde{\mathfrak{g}})$, avec
peut-être un tilde sur le premier $\mathbb{P}$ ; la page 52, (21), a les tildes
dans l'ordre qu'on suit ici, conforme à la convention des pages.}. Pour un tel
sous-groupe $\mathcal{T}$, la page note : $\mathcal{T}$ est son propre centralisateur ;
$X - Y$ ($Y$ le diviseur de degré $2$ correspondant) est un torseur sous
$\mathcal{T}$ ; $N(\mathcal{T})/\mathcal{T}$ s'injecte dans $\underline{\mathrm{Aut}}(\mathcal{T})$ et
contient $-\mathrm{id}_{\mathcal{T}}$\footnote{Les quatre lignes, la note de droite et la
marge sont en grande partie illisibles ; on ne garde que ce qui se lit. Sur
« son propre centralisateur », voir les pages 40 et 55 : en caractéristique
$2$ c'est faux pour le centralisateur de l'algèbre de Lie.}.

\pagerange{36}{40}
\subsection*{Commutant d'un endomorphisme en rang $2$ (pages 36 à 40)}

Soient $V$ un Module de rang $2$ sur $S$ et $u \in \underline{\mathrm{End}}(V)$ nulle part
scalaire : $u_s \notin k(s)\cdot 1$ pour tout $s \in S$. Le commutant de $u$ dans
$\underline{\mathrm{End}}(V)$ est le sous-fibré $\mathcal{O}_S + \mathcal{O}_S u$ : c'est le noyau
de $v \mapsto [u, v]$, dont le rang est $2$ en tout point (page 38), et il
contient $\mathcal{O}_S + \mathcal{O}_S u$, qui est de rang $2$. Le commutant dans
$\mathrm{GL}(V)$ est formé des $\alpha + \beta u$ avec
\[
  (2) \qquad \det(\alpha + \beta u) = \alpha^2 + \alpha\beta\, \mathrm{Tr}\, u + \beta^2 \det u
\]
inversible, et son image $\mathcal{T}$ dans $\mathrm{GP}(V) = \mathrm{PGL}(V)$ est
\[
  \mathcal{T} = C_{\mathrm{PGL}(V)}(u) \simeq \mathbb{P}^1 - Q_{\mathrm{Tr}\,u,\ \det u} ,
\]
où $Q_{\tau, \delta} \subset \mathbb{P}^1$ est le diviseur de degré $2$ d'équation
$\alpha^2 + \tau\alpha\beta + \delta\beta^2 = 0$, étale si et seulement si
$\tau^2 - 4\delta$ est inversible, c'est-à-dire si les valeurs propres de $u$ sont
distinctes en chaque point. $\mathcal{T}$ est un tore si ce diviseur est étale, une
forme de $\mathbb{G}_a$ en un point où il est double\footnote{Le début de la
phrase de la page 36 sur $Q_{\tau,\delta}$ est illisible ; on donne le sens attendu.
La dernière précision est de l'édition.}.

$\mathcal{T}$ ne dépend que de $\mathcal{O}_S + \mathcal{O}_S u$, ou de son image
$\mathfrak{h}$ dans $\underline{\mathrm{End}}(V)/\mathcal{O}_S = \mathrm{Lie}\,\mathrm{PGL}(V)$,
sous-fibré de rang $1$ qui est $\mathrm{Lie}\,\mathcal{T}$. D'où une correspondance
entre sous-fibrés $\mathfrak{h}$ de rang $1$ de $\mathfrak{g} = \mathrm{Lie}\,\mathrm{PGL}(V)$
et sous-groupes commutatifs lisses de dimension relative $1$ de
$\mathrm{PGL}(V)$\footnote{La page écrit « $\mathfrak{g} = \underline{\mathrm{GP}}(V) = G$ »,
là où l'on attend $\mathrm{Lie}$. Elle nomme $H$ le sous-groupe ; voir les
conventions.}.

\textbf{$\mathcal{T}$ est-il le centralisateur de $\mathfrak{h}$ ?} Deux questions se
mêlent sur les pages 37, 39 et 40.

\emph{Le centralisateur d'un élément $h$ de $\mathcal{T}$ (pages 37 et 39).} Si
$g h g^{-1} = \lambda h$ dans $\mathrm{GL}(V)$, le déterminant donne $\lambda^2 = 1$,
donc $g h^2 g^{-1} = h^2$ et $g$ commute à $h^2$ ; si $h^2$ n'est pas scalaire,
c'est-à-dire si $h$ n'est pas une « réflexion », $g \in \mathcal{T}$. Mais en
caractéristique $2$ il peut arriver que $h^2 = 1$ pour tout $h \in \mathcal{T}$ (le cas
unipotent : $h = 1 + tv$, $v^2 = 0$). La page passe alors à la section
universelle $\alpha + \beta u_0$, $\alpha, \beta$ indéterminées : si
$g u_0 g^{-1} = \lambda_0 u_0$ et $g(\alpha + \beta u_0)g^{-1} = \lambda(\alpha,
\beta)(\alpha + \beta u_0)$, la comparaison des coefficients de $1$ donne
$\lambda(\alpha, \beta)\alpha = \alpha$ dans l'anneau localisé
$\mathcal{O}_S[\alpha, \beta]_{\det(\alpha + \beta u_0)}$, d'où
$\lambda = 1$\footnote{La page écrit « $\alpha\lambda(\alpha, \beta) \equiv 1$ » ;
c'est $\lambda\alpha = \alpha$ que donne la comparaison, et la conclusion $\lambda = 1$
est celle de la page.}.

\emph{Le centralisateur de $\mathfrak{h}$ (page 40).} Si $g$ centralise $\mathfrak{h}$,
alors $gug^{-1} = u + \lambda$ ; la trace donne $2\lambda = 0$. Si $2$ est inversible,
$\lambda = 0$ et $g \in \mathcal{T}$ : « gagné ». En caractéristique $2$, non : le
déterminant donne $\lambda(\mathrm{Tr}\,u + \lambda) = 0$, et $\lambda = \mathrm{Tr}\,u$
convient. Par exemple $u = \mathrm{diag}(\zeta, \zeta + 1)$ avec
$\zeta^2 + \zeta + 1 = 0$ est conjugué à $u + 1 = \mathrm{diag}(\zeta + 1, \zeta)$ sans
lui être égal : le centralisateur de $\mathfrak{h}$ contient l'élément qui échange
les deux droites propres, et contient strictement $\mathcal{T}$, tout en restant
lisse de dimension relative $1$ (c'est le normalisateur du tore)\footnote{Ce
paragraphe de la page 40 est barré sur la page, mais l'exemple est juste et la
page le reprend aussitôt sous forme positive (« la translation par
$\zeta + \zeta'$ le transforme en $(\zeta', \zeta)$, qui lui est conjugué ») avant de
finir sur un point d'interrogation. La parenthèse finale est de l'édition.}.
La même chose se lit dans l'algèbre de Lie : en caractéristique $2$, dans
$\mathfrak{gp}(1)$ où $[H, X] = X$, $[H, Y] = Y$, $[X, Y] = 0$, le centralisateur de
$X$ est $\{h = 0\}$, de dimension $2$, « trop grande », celui de $H$ est
$\mathcal{O}H$.

\pagerange{41}{44}
\subsection*{Sous-groupes d'ordre premier impair (pages 41 à 44, et le haut de la page 38)}

La page 41 commence au milieu d'une démonstration dont l'énoncé est à la page
42, et le haut de la page 38 semble en être la fin\footnote{Rattachement proposé
par la transcription (note de la page 38), non marqué par l'auteur.}.

\textbf{Énoncé (page 42).} \emph{Soient $G$ une forme de $\mathrm{PGL}_2$ sur $S$, $p$ un
nombre premier impair, $\mu \subset G$ un sous-groupe fini étale de rang $p$.}
\begin{itemize}
  \item[a)] \emph{Le centralisateur $\mathcal{T}$ de $\mu$ est lisse, commutatif, de
    dimension relative $1$.}
  \item[b)] \emph{Si $p$ est inversible sur $S$, soit $N$ le normalisateur de $\mu$ ;
    alors $\mathcal{T} = N^{\circ}$, et $N/\mathcal{T} \simeq \{\pm 1\}_S$ opère sur $\mathcal{T}$
    par $h \mapsto h^{-1}$.}
\end{itemize}
\footnote{L'hypothèse « $p$ inversible » dans b) est de l'édition. Sur une fibre
de caractéristique $p$, $\mu$ est engendré par un unipotent $u = 1 + v$, le
centralisateur est $\mathbb{G}_a$, et le normalisateur contient
$\mathrm{diag}(a, 1)$ pour tout $a \in \mathbf{F}_p^{\times}$, qui envoie $u$ sur $u^a$ :
$N/\mathcal{T} \simeq \mathbf{F}_p^{\times}$ tout entier, égal à $\{\pm 1\}$ seulement pour
$p = 3$. La page traite bien le cas $p = \mathrm{car}\,k$ à la page 41, mais
seulement pour le commutant. Quant à la connexité des fibres de $\mathcal{T}$, que la
page restreint de façon peu lisible aux caractéristiques résiduelles $\neq p$,
elle vaut partout : tore maximal en caractéristique $\neq p$, $\mathbb{G}_a$ en
caractéristique $p$. Un c), sur une section de $N/\mathcal{T}$ d'image $-1$ dans
$(\mathbf{Z}/p\mathbf{Z})^{\times}$, est barré.}

\textbf{Démonstration (pages 42, 41, 38).} Localement $G = \widetilde{G}/\mu_2$ avec
$\widetilde{G} = \mathrm{SL}(V)$ ; comme $p$ est impair, $\mu$ se relève de façon
unique en $\widetilde{\mu} \simeq \mu$ dans $\widetilde{G}$, et il suffit de voir que
l'image du centralisateur de $\widetilde{\mu}$ est le centralisateur de $\mu$. Sur un
corps $k$ algébriquement clos de caractéristique $\neq p$, un générateur de
$\widetilde{\mu}$ a des valeurs propres $\zeta, \zeta^{-1}$, $\zeta$ racine primitive
$p$-ième, distinctes puisque $p$ est impair ; le commutant est le tore diagonal,
le normalisateur est formé des matrices diagonales et de leurs produits par une
« symétrie », d'où $N/\mathcal{T} \simeq \{\pm 1\} \subset (\mathbf{Z}/p\mathbf{Z})^{\times}$. En
caractéristique $p$, $u = 1 + v$, $v^2 = 0$, $v \neq 0$, et le commutant dans
$\mathrm{SL}$ est formé des $\alpha + \beta v$ avec $\alpha^2 = 1$, $\alpha = \pm 1$ ;
son image dans $\mathrm{PGL}$ est $\mathbb{G}_a$, lisse et connexe. Dans le cas général,
l'image inverse par $N \to \underline{\mathrm{Aut}}(\mu) = (\mathbf{Z}/p\mathbf{Z})^{\times}_S$ de la
section unité est un sous-groupe ouvert de $N$ qui coïncide fibre à fibre avec
$N^{\circ}$, donc lui est égal ; et $N/N^{\circ}$ est un sous-groupe étale ouvert de
$(\mathbf{Z}/p\mathbf{Z})^{\times}_S$ qui coïncide fibre à fibre avec $\{\pm 1\}$, donc lui est
égal (haut de la page 38) — sous l'hypothèse de b)\footnote{La lissité est dite
« connue pour tout sous-groupe de type multiplicatif dans un groupe lisse »
lorsque $p$ est premier à la caractéristique ; la page ne va pas plus loin, et
le cas de caractéristique résiduelle $p$ sur une base n'est pas traité.}.

\textbf{Une question (page 43).} Tout sous-groupe lisse $\mathcal{T}$ de dimension
relative $1$ de $G$ est-il de la forme $C(\mathfrak{h})$, déterminé par son algèbre de
Lie ? Oui s'il est commutatif et si les caractéristiques résiduelles sont
$\neq 2$, car il est contenu dans $C(\mathfrak{h})$, qui a les mêmes fibres. Mais la
commutativité est-elle automatique ? Sur un corps oui — tout groupe lisse connexe
de dimension $1$ est commutatif — ; sur une base, c'est une question de
déformation de sous-groupes, « liée à un $H^2$ de Hochschild ». Le normalisateur
de $\mathcal{T}$, qui est aussi celui de $\mathfrak{h}$, n'est pas en général de
dimension relative $1$ : sur un corps, quand $\mathcal{T}$ est unipotent, c'est un
Borel. La phrase s'arrête au haut de la page 44 sur « une structure décrite par
\ldots{} rigidifiée »\footnote{Lecture incertaine ; le reste de la page 44 est
blanc, la page 45 aussi.}.

\pagerange{46}{55}
\subsection*{Algèbres de quaternions (pages 46 à 55)}

\textbf{Le programme (page 46).} À une algèbre de quaternions $\mathcal{M}$ sur $S$ — une
algèbre d'Azumaya de rang $4$ — on associe une forme $G$ de $\mathrm{PGL}_2$, une
forme $\widetilde{G}$ de $\mathrm{SL}_2$, leurs algèbres de Lie $\mathfrak{g}$ et
$\widetilde{\mathfrak{g}}$, un module spécial quadratique $(E, q, \omega)$ de rang $3$
avec $E = \widetilde{\mathfrak{g}}$, une conique lisse $X$ dans un fibré en plans
projectifs $P = \check{\mathbb{P}}(E)$, $X = \mathrm{Var}(q)$, et $X$ seule, fibré en
droites projectives\footnote{La page lit « rg 2 » pour le module spécial
quadratique, qui est de rang $3$ ; le modèle de $\widetilde{\mathfrak{g}}$ dans la liste
est illisible, et l'on attend $\mathfrak{sl}(2)$ (page 59).}. Les foncteurs
$\mathcal{M} \mapsto G, \widetilde{G}, (E, q, \omega), (P \supset X), X$ sont des
équivalences de catégories fibrées, et $\mathcal{M} \mapsto \mathfrak{g},
\widetilde{\mathfrak{g}}$ en sont au-dessus de $\operatorname{Spec}\mathbf{Z}[\tfrac{1}{2}]$.
C'est affirmé, non démontré ; c'est aujourd'hui le dictionnaire classique entre
algèbres de quaternions, formes quadratiques ternaires (par l'algèbre de Clifford
paire) et coniques\footnote{Références de l'édition, citées de mémoire : M.-A.
Knus, \emph{Quadratic and Hermitian Forms over Rings} (1991), pour l'algèbre de
Clifford paire sur une base ; et, pour les ordres de quaternions et les formes
ternaires sur un anneau, des travaux postérieurs (Gross--Lucianovic, 2009 ;
J. Voight, \emph{Quaternion Algebras}, 2021). Rien sur les pages ne cite de
littérature. Le dossier 77 (page 89) part du même dictionnaire : un module
spécial quadratique de rang $3$ y est identifié aux éléments de trace réduite
nulle d'une algèbre de quaternions tordue, avec $q = -\det$.}.

On se place dans le cas trivial $\mathcal{M} = V^{\vee} \otimes V \simeq
\underline{\mathrm{End}}(V)$, $V$ de rang $2$ (1).

\textbf{La conique (page 46).} À un sous-fibré en droites $L \subset V$ on associe
$L^{\perp} \otimes L \subset V^{\vee} \otimes V = \mathcal{M}$ : ce sont les endomorphismes
d'image dans $L$ et nuls sur $L$, de trace et de déterminant nuls, et
$L^{\perp} \otimes L$ est une droite isotrope de $\widetilde{\mathfrak{g}}$. D'où un
isomorphisme canonique
\[
  (2) \qquad \check{\mathbb{P}}(V) \xrightarrow{\ \sim\ } X , \qquad L \longmapsto L^{\perp} \otimes L ,
\]
où $\check{\mathbb{P}}(V) \simeq \mathbb{P}(V)$, puisque $V$ est de rang $2$\footnote{Le
dossier 89 (pages 15 à 18) reconstitue un module libre de rang $2$ à partir de
ses droites et de son déterminant ; c'est la même géométrie de la droite
projective, vue de l'autre côté.}.

\textbf{La quadrique des endomorphismes de rang $1$ (pages 47 et 48).} Soient
$Q_{\mathcal{M}}(u) = -\det u$, forme quadratique non dégénérée sur $\mathcal{M}$, et
$\mathcal{Q} = \mathrm{Var}(Q_{\mathcal{M}}) \subset \check{\mathbb{P}}(\mathcal{M})$\footnote{La page note
$\Delta$ cette forme et $Q$ la quadrique ; on renomme pour ne pas confondre avec
le $\Delta$ de la première suite.}. Les points de $\mathcal{Q}$ sont les droites
$\mathcal{O}_S u$ engendrées par un $u$ nulle part nul de déterminant nul, c'est-à-dire
strictement de rang $1$ : $u$ se factorise en $V \twoheadrightarrow I \hookrightarrow V$,
$I$ inversible, et se donner $\mathcal{O}_S u$ revient à se donner indépendamment
$\mathrm{Ker}\,u$ et $\mathrm{Im}\,u$. D'où l'isomorphisme
\[
  (6)\text{--}(7) \qquad \mathcal{Q} \xrightarrow{\ \sim\ } \check{\mathbb{P}}(V) \times \check{\mathbb{P}}(V)
  \simeq X \times X , \qquad u \longmapsto (\mathrm{Ker}\,u, \mathrm{Im}\,u) ,
\]
d'inverse $(L, M) \mapsto L^{\perp} \otimes M$ : c'est le plongement de Segre de
$\mathbb{P}^1 \times \mathbb{P}^1$ comme quadrique de $\mathbb{P}^3$. Les droites
$\underline{L} = L^{\perp} \otimes L$, $\underline{M} = M^{\perp} \otimes M$ sont dans
$\widetilde{\mathfrak{g}}$ et isotropes, et la droite $J = \mathcal{O}_S u$ vérifie
$J \circ \underline{L} = 0$ et $\underline{M} \circ J = 0$ (8)--(9)\footnote{Le
premier $\mathfrak{g}$ de (8) porte une barre sur la page ; c'est
$\widetilde{\mathfrak{g}}$ dans les deux lignes.}.

\textbf{Algèbres spéciales (page 48).} Pour $J \in \mathcal{Q}(S)$, de générateur local
$u$, on a $u^2 = (\mathrm{Tr}\,u)\,u$, donc $\mathcal{A} = \mathcal{O}_S \oplus J$ est une
sous-algèbre de rang $2$ de $\mathcal{M}$, contenant $1$ — une sous-algèbre
\emph{spéciale} —, et $J$ en est un idéal de quotient $\mathcal{O}_S$. Se donner $J$
revient à se donner $\mathcal{A}$ et un homomorphisme d'algèbres
$\rho : \mathcal{A} \to \mathcal{O}_S$, celui de noyau $J$. D'où
\[
  (12) \qquad \mathcal{Q} \simeq \bigl\{ (\mathcal{A}, \rho) : \mathcal{A} \subset \mathcal{M} \ \text{sous-algèbre spéciale},\
  \rho : \mathcal{A} \to \mathcal{O}_S \ \text{homomorphisme d'algèbres} \bigr\}
\]
\footnote{Réciproque, de l'édition : si $\mathcal{A} = \mathcal{O}_S[u]$ et
$\rho(u) = c$, le noyau est engendré par $u - c$, et
$\det(u - c) = c^2 - c\,\mathrm{Tr}\,u + \det u = 0$ par Cayley--Hamilton.}.

\textbf{L'involution canonique (page 49).} La symétrie des deux facteurs de
$X \times X$ est induite par un anti-automorphisme canonique de $\mathcal{M}$ : le choix
d'une base de $\det V$ donne $\alpha : V \simeq V^{\vee}$, et
$u' = \alpha^{-1}\,{}^{t}u\,\alpha$ ne dépend pas de ce choix. On a $(uv)' = v'u'$ et
\[
  (16)\text{--}(17) \qquad u' = (\mathrm{Tr}\,u)\cdot 1 - u , \qquad u + u' = \mathrm{Tr}\,u , \qquad uu' = \det u :
\]
c'est la conjugaison des quaternions, l'adjugée en termes de matrices. Elle
échange noyau et image des $u$ de rang $1$ ($Q_{\mathcal{M}}(u) = 0$, $u$ nulle part
nul)\footnote{L'encre donne « $\Gamma(u) = 0$ » ; c'est la forme de (5).}. En termes
de (12), elle envoie $(\mathcal{A}, \rho)$ sur $(\mathcal{A}, \rho')$, $\rho'(u) = \rho(u')$ :
$\mathcal{A}' = \mathcal{A}$, et sur $\mathcal{A}$ commutative l'involution est un
automorphisme, l'automorphisme non trivial de l'algèbre de rang $2$.

\textbf{Le schéma des algèbres spéciales (page 50).} Une algèbre spéciale
$\mathcal{A}$ est déterminée par la droite $\mathcal{A}/\mathcal{O}_S$ de $\mathfrak{g} = \mathcal{M}/\mathcal{O}_S$ ;
le schéma $\Sigma$ des algèbres spéciales est donc $\check{\mathbb{P}}(\mathfrak{g}) =
\mathbb{P}(\widetilde{\mathfrak{g}})$, et $(\mathcal{A}, \rho) \mapsto \mathcal{A}$ est un épimorphisme
$\mathcal{Q} \to \Sigma$ qui identifie $\mathcal{Q}$ au spectre de l'algèbre spéciale universelle
$\mathcal{A}_{\Sigma}$ (20). La page esquisse ensuite, puis barre, une liste de
conditions équivalentes, reprise à la page suivante.

\textbf{Proposition (page 51).} \emph{Soient $\mathcal{A}$ une sous-algèbre spéciale,
$L \subset V$ un sous-fibré en droites, $\underline{L} = L^{\perp} \otimes L$. Sont
équivalentes :}
\begin{itemize}
  \item[a)] \emph{$\underline{L} \subset \mathcal{A}^{\perp}$ (orthogonal pour
    $\mathrm{Tr}(uv)$) ;}
  \item[b)] \emph{$\mathcal{A}$ stabilise $L$ ;}
  \item[c)] \emph{il existe un idéal $J$ de $\mathcal{A}$, sous-fibré de rang $1$, avec
    $J \cdot L = 0$ ;}
  \item[d)] \emph{il existe un sous-fibré en droites $M \subset V$ avec
    $L^{\perp} \otimes M \subset \mathcal{A}$.}
\end{itemize}
\emph{Alors $J$ et $M$ sont uniques, $J = L^{\perp} \otimes M$ est le noyau de
l'homomorphisme $\mathcal{A} \to \underline{\mathrm{End}}(L) = \mathcal{O}_S$ déduit de b), et,
pour $\mathcal{A}$ fixée, $L \mapsto (\mathcal{A} \to \mathcal{O}_S)$ est une bijection entre les $L$
vérifiant a) et les homomorphismes d'algèbres $\mathcal{A} \to \mathcal{O}_S$.}

\emph{Démonstration (pages 51 et 52).} Si $L = \mathcal{O}x$ et $L^{\perp} = \mathcal{O}x'$,
$\mathrm{Tr}(u \cdot x' \otimes x) = \langle ux, x' \rangle$, nul si et seulement si
$u(L) \subset L$ : a) $\Leftrightarrow$ b). Le noyau $J$ de
$\mathcal{A} \to \underline{\mathrm{End}}(L)$ annule $L$ : b) $\Rightarrow$ c) ; réciproquement,
$\mathcal{A}/J \simeq \mathcal{O}_S$, donc $\mathcal{A} = \mathcal{O}_S + J$ stabilise $L$, et $J$ est
déterminé comme ce noyau. Sous ces conditions, un générateur local $u$ de $J$
est de rang $1$ avec $\mathrm{Ker}\,u = L$ ; posant $M = \mathrm{Im}\,u$, on a
$L^{\perp} \otimes M = J \subset \mathcal{A}$ : d). Réciproquement, si
$L^{\perp} \otimes M \subset \mathcal{A}$, c'est un idéal de rang $1$ qui annule $L$, donc
$J$, et $M = \mathrm{Im}\,u$ est déterminé ; et d) $\Rightarrow$ b) directement.

\textbf{Corollaire (pages 52 et 53).} Se donner une sous-algèbre spéciale
$\mathcal{A}$ — un sous-fibré de rang $2$ contenant $\mathcal{O}_S\cdot 1$, nécessairement
commutatif — revient à se donner son orthogonal
$\mathcal{A}^{\perp} = \mathcal{H} \subset \widetilde{\mathfrak{g}}$, hyperplan de $\widetilde{\mathfrak{g}}$ :
\[
  (21) \qquad \text{AlgSpéc}(\mathcal{M}) \simeq \mathbb{P}(\widetilde{\mathfrak{g}}) \simeq
  \check{\mathbb{P}}(\mathfrak{g}) .
\]
Comme $P = \check{\mathbb{P}}(\widetilde{\mathfrak{g}})$, un hyperplan $\mathcal{H}$ de
$\widetilde{\mathfrak{g}}$ est une droite $D$ du plan $P$, et
\[
  (22) \qquad \check{P} = \mathrm{Dr}(P) \xrightarrow{\ \sim\ } \mathrm{Div}_2^{+}(X), \qquad D \longmapsto D \cap X = Y ,
\]
est un isomorphisme : une droite du plan coupe la conique lisse en un diviseur
de degré $2$, et tout diviseur de degré $2$ s'obtient ainsi d'une seule
façon\footnote{Vrai en toute caractéristique : le plongement de $X \simeq \mathbb{P}^1$
dans $P$ est le plongement par $\mathcal{O}(2)$, et les diviseurs de degré $2$ sont les
sections de $\mathcal{O}(2)$ à scalaire près.}. Le diviseur $Y$ est formé des droites
isotropes $\underline{L}$ contenues dans $\mathcal{H}$, et l'on trouve un isomorphisme
canonique
\[
  (23) \qquad Y \simeq \operatorname{Spec} \mathcal{A} .
\]
Dans $X \simeq \check{\mathbb{P}}(V)$, $Y$ est formé des $L$ stables par $\mathcal{A}$, ou, si
$\mathcal{A} = \mathcal{O}_S + \mathcal{O}_S u$, par $u$ : si $u$ est inversible, $Y = X^{u}$, le
sous-schéma des points fixes de $u$ (24).

\textbf{Le stabilisateur, et les sous-groupes jacobiens (pages 53 à 55).}
$G = \mathcal{M}^{\times}/\mathbb{G}_m$ opère sur tout cela. \emph{Proposition (page 54) :
$g \in G(S)$ stabilise $Y$ si et seulement si $g\mathcal{A} g^{-1} = \mathcal{A}$ ; alors son
action sur $Y$ est, via (23), son action sur $\operatorname{Spec}\mathcal{A}$, triviale si et
seulement si $g$ est (localement l'image d'un élément) de $\mathcal{A}^{\times}$.} À
$\mathcal{A}$ on associe
\[
  (25) \qquad \mathcal{T} = \mathcal{A}^{\times}/\mathbb{G}_m \subset G ,
\]
sous-groupe commutatif, lisse, à fibres connexes, de dimension relative $1$ — un
tore si $\mathcal{A}$ est étale, $\mathbb{G}_a$ là où $\mathcal{A}$ est
$\mathcal{O}[\varepsilon]/(\varepsilon^2)$. La correspondance $\mathcal{A} \mapsto \mathcal{T}$ est
injective, $\mathcal{A}$ étant engendrée par les sections de $\mathcal{A}^{\times}$, image
inverse de $\mathcal{T}$ dans $\mathcal{M}^{\times}$. La page appelle ces sous-groupes
« \emph{jacobiens} »\footnote{Le mot est de la page. Une explication possible, qui
est de l'édition : $\mathcal{A}^{\times}/\mathbb{G}_m$ est la jacobienne généralisée de
Rosenlicht de la droite projective $X$ pour le module $Y$ de degré $2$ — $\mathbb{G}_m$
si $Y$ est formé de deux points, $\mathbb{G}_a$ si $Y$ est un point double (J.-P.
Serre, \emph{Groupes algébriques et corps de classes}, 1959).} : ce sont les
$\mathcal{T}$ des pages 36 et 37.

Enfin (page 55), si $\mathfrak{h} = \mathrm{Lie}\,\mathcal{T}$ et si $\widetilde{\mathcal{T}}$ est le
centralisateur de $\mathfrak{h}$ dans $G$, alors $\mathcal{T} \subset \widetilde{\mathcal{T}}$ puisque
$\mathcal{T}$ est commutatif ; la page conclut à l'égalité, $\widetilde{\mathcal{T}}$ étant
connexe et les deux lisses de dimension relative $1$. C'est juste lorsque $2$ est
inversible sur $S$ ; en caractéristique $2$, la page 40 a montré que le
centralisateur de $\mathfrak{h}$ peut contenir strictement $\mathcal{T}$ et n'être pas
connexe\footnote{La page identifie ici « le sous-groupe jacobien que
$\mathfrak{h}$ définit » au centralisateur de $\mathfrak{h}$. Une note marginale
oblique, « cor. p. 52 » (lecture incertaine), signale peut-être une correction à
faire.}.

\section*{III. Dictionnaire lié à $\mathrm{GP}(1) \simeq \mathrm{SO}(3)$ (pages 56 à 78)}

Le titre est inscrit sur la couverture rose de la page 56.

\pagerange{57}{59}
\subsection*{Les neuf catégories (pages 57 à 59)}

Soit $S$ un schéma. On va définir des équivalences entre les catégories
suivantes, qui, pour $S$ variable, sont des catégories fibrées sur
$\mathrm{Sch}_{/S}$, en fait des \emph{gerbes} pour la topologie fppf, formant « un
système transitif d'équivalences » :
\begin{itemize}
  \item[$C_0$] $\mathrm{Tors}(\mathrm{PGL}_{2,S})$, les torseurs à droite sous $\mathrm{PGL}_2$ ;
  \item[$C_1$] $\mathrm{Fo}(\mathrm{PGL}_{2,S})$, les formes $G$ de $\mathrm{PGL}_2$ ;
  \item[$C'_1$] $\mathrm{Fo}(\mathrm{SL}_{2,S})$, les formes $\widetilde{G}$ de $\mathrm{SL}_2$ ;
  \item[$C_2$] $(*)$ $\mathrm{Fo}(\mathfrak{pgl}_{2,S})$, les formes $\mathfrak{g}$ de
    l'algèbre de Lie $\mathfrak{gp}(1)$ ;
  \item[$C'_2$] $(*)$ $\mathrm{Fo}(\mathfrak{sl}_{2,S})$, les formes $\widetilde{\mathfrak{g}}$ de
    $\mathfrak{sl}(2)$ ;
  \item[$C_3$] $\mathrm{Quat}(S)$, les algèbres d'Azumaya $\mathcal{M}$ de rang $4$ ;
  \item[$C_4$] $\mathrm{Modspquad}(S)$, les modules spéciaux quadratiques
    $(E, q, \omega)$ de rang $3$ ;
  \item[$C_5$] $\mathrm{Drpr}(S)$, les fibrés en droites projectives $X$ ;
  \item[$C'_5$] $\mathrm{Quad}(S)$, les couples $(P, X)$ d'un fibré en plans projectifs
    et d'une conique lisse $X \subset P$.
\end{itemize}
\footnote{Deux lignes, $C''_1 = \mathrm{Fo}(\mathrm{GL}(2))$ et
$C''_2 = \mathrm{Fo}(\mathfrak{gl}(2))$ (avec $(**)$), sont barrées ; elles expliquent
que la page 59 parle de « 11 catégories » alors qu'il en reste neuf, et que la
page 58 parle de « 3 » catégories marquées $(*)$ et de « 6 autres » (lecture
incertaine), là où il en reste deux et sept. La liste est disposée en deux
colonnes avec des accolades ; on la donne dans l'ordre de la page.}

Le signe $(*)$ marque les deux catégories dont l'équivalence avec les autres ne
vaut que si $2$ est inversible sur $S$ ; en général on n'a que des foncteurs des
autres vers elles, qui ne sont pas pleinement fidèles — ce que la dimension $6$
de la page 27 rendait inévitable. Un dictionnaire complet comprendrait la
description explicite des $9 \times 8 = 72$ équivalences et des $9 \cdot 8 \cdot 7 =
504$ isomorphismes de transitivité. « On fera un peu moins, mais on essaiera de
donner une image claire de la situation. »

\textbf{Principe général.} Sur $S_0 = \operatorname{Spec}\mathbf{Z}$, on choisit dans chaque
catégorie un objet typique $\mathfrak{X}_0$ et une opération de
$G_0 = \mathrm{PGL}_2$ sur lui, et l'on établit que $G_0 \to
\underline{\mathrm{Aut}}(\mathfrak{X}_0)$ est un isomorphisme — sauf dans les deux cas $(*)$,
où l'opération est fidèle mais l'homomorphisme n'est pas un épimorphisme, et
où l'on se borne à $\operatorname{Spec}\mathbf{Z}[\tfrac{1}{2}]$. Comme les catégories
envisagées sont des gerbes, chacune est alors canoniquement équivalente à celle
des torseurs sous $\underline{\mathrm{Aut}}(\mathfrak{X}_0)$, donc sous $G_0$ : c'est
l'\emph{existence} d'un « multi-dictionnaire », qui ne dispense pas de le
développer\footnote{En termes d'aujourd'hui : chacune de ces gerbes est
neutre, de lien $\mathrm{PGL}_2$, et les objets à isomorphisme près sur $S$ sont
classés par $H^1(S, \mathrm{PGL}_2)$ — c'est le langage de la cohomologie non
abélienne de J. Giraud (1971).}.

\pagerange{59}{62}
\subsection*{Représentants typiques et opérations de $G_0$ (pages 59 à 62)}

\textbf{I. Représentants typiques.} Sur $S_0 = \operatorname{Spec}\mathbf{Z}$ : le torseur
trivial $G_0$ ; $G_0 = \mathrm{PGL}_2$ et $\widetilde{G}_0 = \mathrm{SL}_2$ ;
$\mathfrak{g}_0 = \mathfrak{pgl}_2$ et $\widetilde{\mathfrak{g}}_0 = \mathfrak{sl}_2$ ;
$\mathcal{M}_0 = \mathrm{M}_2(\mathcal{O})$ ; $E_0 = \mathcal{O}^3$ avec $q_0 = xy + z^2$ et
$\omega = e_1 \wedge e_2 \wedge e_3$ ; $X_0 = \mathbb{P}^1$ ; $P_0 = \mathbb{P}^2 =
\check{\mathbb{P}}(E_0)$ avec la conique $C_0 = \mathrm{Var}(q_0)$.

\textbf{II. Opérations de $G_0$.} 1) Les translations (la ligne s'arrête sur
« gauche… »). 2) $G_0$ opère sur lui-même par automorphismes intérieurs ; ceux de
$\widetilde{G}_0$ passent au quotient par le centre $\mu_2$, et
$\mathrm{SL}_2/\mu_2 \simeq \mathrm{PGL}_2$ (1) donne une action de $G_0$ sur
$\widetilde{G}_0$. 3) D'où des actions sur $\mathfrak{g}_0$ et $\widetilde{\mathfrak{g}}_0$.
4) L'action de $\mathrm{GL}_2 = \mathcal{M}_0^{\times}$ sur $\mathcal{M}_0$ par conjugaison est
triviale sur le centre $\mathbb{G}_m$ et passe au quotient
$\mathrm{GL}_2/\mathbb{G}_m = \mathrm{PGL}_2$ (2).

5) (page 61) La forme $\xi \mapsto -\det\xi$ sur $\mathcal{M}_0$, restreinte aux matrices
de trace nulle
\[
  (3)\text{--}(4) \qquad \xi = \begin{pmatrix} z & x \\ y & -z \end{pmatrix}, \qquad
  q_0(\xi) = -\det\xi = xy + z^2 ,
\]
est invariante par conjugaison, et l'identification
\[
  (6) \qquad E_0 = \mathcal{O}^3 \xrightarrow{\ \sim\ } \widetilde{\mathfrak{g}}_0 = \mathfrak{sl}_2 , \qquad
  (x, y, z) \longmapsto \begin{pmatrix} z & x \\ y & -z \end{pmatrix}
\]
donne une action de $G_0$ sur $(E_0, q_0)$. Elle préserve aussi $\omega$, car $G_0$
opère trivialement sur tout Module inversible, $\underline{\mathrm{Hom}}(\mathrm{PGL}_2,
\mathbb{G}_m) = 0$ (page 62). 6) $\mathrm{GL}_2$ opère sur $\mathbb{P}^1$, $\mathbb{G}_m$
trivialement, d'où l'action de $G_0$ ; et l'action sur $(E_0, q_0)$ en donne une
sur $(P_0, C_0)$.

« Nous admettrons (pour ne pas nous éterniser…) que tous les homomorphismes
\[
  (7) \qquad G_0 = \mathrm{PGL}_2 \longrightarrow \underline{\mathrm{Aut}}(\mathfrak{X}_0)
\]
sont des isomorphismes », en dehors des cas $(*)$, et en marge : « Expliciter le
théorème d'isomorphie »\footnote{Tous ces énoncés sont connus aujourd'hui sur
$\mathbf{Z}$ : $\underline{\mathrm{Aut}}(\mathbb{P}^1) = \mathrm{PGL}_2$,
$\underline{\mathrm{Aut}}(\mathrm{M}_2) = \mathrm{PGL}_2$ (Skolem--Noether), $\mathrm{SO}(xy + z^2)
\simeq \mathrm{PGL}_2$, $\underline{\mathrm{Aut}}(\mathrm{SL}_2) = \underline{\mathrm{Aut}}(\mathrm{PGL}_2) =
\mathrm{PGL}_2$ (pas d'automorphisme extérieur pour le type $A_1$), et
$\underline{\mathrm{Aut}}(\mathbb{P}^2, C_0) \simeq \underline{\mathrm{Aut}}(C_0)$, la conique étant
plongée par un système linéaire intrinsèque.}.

\pagerange{63}{65}
\subsection*{Les foncteurs : torsion, torseur des isomorphismes, automorphismes (pages 63 à 65)}

\textbf{III. De $C_0$ vers les autres.} C'est la \emph{torsion}
\[
  (8) \qquad R \longmapsto R \overset{\mathrm{PGL}_{2,S}}{\wedge} \mathfrak{X}_{0,S} ,
\]
et le théorème d'isomorphie (7) dit que ces foncteurs sont des équivalences (sur
$\mathbf{Z}[\tfrac{1}{2}]$ pour les cas $(*)$). Vers $C_5$ :
\[
  (9) \qquad R \longmapsto R/B_{0,S} ,
\]
où $B_0 \subset G_0$ est le Borel type, image dans $\mathrm{PGL}_2$ du groupe
$\widetilde{B}_0$ des matrices $\begin{pmatrix} \lambda & \mu \\ 0 & \lambda^{-1} \end{pmatrix}$
de $\mathrm{SL}_2$, stabilisateur de la droite $\mathcal{O} e_1$ de $\mathcal{O}^2$ (10). Le
point correspondant $s_0$ de $X_0 = \mathbb{P}^1$ a pour stabilisateur $B_0$, et
$g \mapsto g\cdot s_0$ induit un isomorphisme $G_0/B_0 \simeq X_0$ compatible aux
translations (13) ; tordant par $R$, on obtient
$R \wedge X_{0,S} \simeq R/B_{0,S}$ (14)\footnote{La page 64 nomme d'abord ce point
« la section nulle $s_0$ », puis « la section $\infty$, correspondant à la droite
$\mathcal{O}_S e_1$ », et lui donne pour stabilisateur (11) les matrices
\emph{triangulaires inférieures}, transposées de (10). Les deux lectures se
concilient selon que $\mathbb{P}(\mathcal{O}^2)$ paramètre des droites (action naturelle :
le stabilisateur de $\mathcal{O} e_1$ est triangulaire supérieur) ou des quotients de
rang $1$, comme dans la convention des pages (action contragrédiente : il est
triangulaire inférieur). Les deux sous-groupes sont conjugués, et (13) vaut avec
l'un ou l'autre.}.

\textbf{Des autres vers $C_0$.} C'est
\[
  (15)\text{--}(17) \qquad \mathfrak{X} \longmapsto \underline{\mathrm{Isom}}(\mathfrak{X}_{0,S}, \mathfrak{X}) ,
\]
torseur à droite sous $\underline{\mathrm{Aut}}(\mathfrak{X}_{0,S}) \simeq \mathrm{PGL}_{2,S}$ par
(7)\footnote{Dans (15), le premier $\mathfrak{X}$ porte un indice $0$ de trop.}.

\textbf{IV. Vers $C_1$.} C'est $\mathfrak{X} \mapsto \underline{\mathrm{Aut}}(\mathfrak{X})$ (18),
avec deux descriptions plus terre à terre : $\widetilde{G} \mapsto
\widetilde{G}/\mathrm{Cent}(\widetilde{G})$ (19), le centre étant toujours isomorphe à
$\mu_2$, de façon unique car $\underline{\mathrm{Aut}}(\mu_2) = 1$ ; et
$\mathcal{M} \mapsto \mathcal{M}^{\times}/\mathbb{G}_m$ (20).

\pagerange{66}{68}
\subsection*{Revêtement, algèbres de Lie, sous-groupes de Borel (pages 66 à 68)}

\textbf{De $C_1$ vers les autres.} Ou bien l'on tord $\mathfrak{X}_{0,S}$ par le
torseur $\underline{\mathrm{Isom}}(G_{0,S}, G)$ (21) — torseur sous
$\underline{\mathrm{Aut}}(G_0) = G_0$, puisque $\mathrm{PGL}_2$ est adjoint —, ou bien
l'on a des descriptions directes :
\begin{itemize}
  \item $G \mapsto \widetilde{G}$, le revêtement simplement connexe, schéma en
    groupes lisse à fibres géométriquement connexes et simplement connexes muni
    d'une isogénie centrale $\widetilde{G} \to G$, unique à isomorphisme unique près
    (22) ;
  \item $G \mapsto \mathfrak{g} = \underline{\mathrm{Lie}}(G)$ et $G \mapsto
    \widetilde{\mathfrak{g}} = \underline{\mathrm{Lie}}(\widetilde{G})$, équivalences
    $\mathrm{Fo}(\mathrm{PGL}_2) \to \mathrm{Fo}(\mathfrak{pgl}_2)$, $\mathrm{Fo}(\mathfrak{sl}_2)$
    au-dessus de $\mathbf{Z}[\tfrac{1}{2}]$ (23)--(24)\footnote{La page surmonte les
    flèches de (23) et (24) du signe $\sim$ sans restriction ; ce sont les cas
    $(*)$, et l'on rétablit la restriction de la page 58.} ;
  \item $G \mapsto (E, q, \omega)$, qui se factorise par $G \mapsto \widetilde{\mathfrak{g}}$
    (« cf plus bas » : c'est (50), page 74) ;
  \item $G \mapsto X = \underline{\mathrm{Bor}}(G)$, le schéma des sous-groupes de Borel
    de $G$, c'est-à-dire des $B$ tels que $(G, B)$ soit localement isomorphe à
    $(G_0, B_0)$ (25) ; si $G$ et $X$ se correspondent,
    $X \xrightarrow{\sim} \underline{\mathrm{Bor}}(G)$ associe à une section locale de $X$ son
    stabilisateur (26).
\end{itemize}

\textbf{V. Vers les formes de $\mathrm{SL}_2$.} Ou bien par $C_0$ (27), ou bien, plus
naturellement, $\mathfrak{X} \mapsto \underline{\mathrm{Aut}}(\mathfrak{X}) \mapsto$ son
revêtement simplement connexe (28) ; depuis $C_3$,
\[
  (29) \qquad \mathcal{M} \longmapsto \widetilde{G} = \mathrm{Ker}\bigl(\mathcal{M}^{\times}
  \xrightarrow{\ \det\ } \mathbb{G}_m\bigr),
\]
$\det$ étant la norme réduite. Inversement, $\widetilde{G} \mapsto \widetilde{\mathfrak{g}}$ (30)
et $\widetilde{G} \mapsto X \simeq \underline{\mathrm{Bor}}(\widetilde{G})$ (31). Pour
comprendre ensemble $G$, $\widetilde{G}$, $\mathfrak{g}$, $\widetilde{\mathfrak{g}}$,
$(E, q, \omega)$, $X$ et $(P, C)$, « c'est $\mathcal{M}$ qui donne peut-être la meilleure
clef » ; on verra en particulier que $\mathfrak{g}$ et $\widetilde{\mathfrak{g}}$ sont
duaux l'un de l'autre, et que $E \simeq \widetilde{\mathfrak{g}}$ canoniquement.

\pagerange{69}{73}
\subsection*{VI. L'algèbre de quaternions comme clef (pages 69 à 73)}

Soit $\mathcal{M}$ une algèbre de quaternions sur $S$, localement (fppf) isomorphe à
$\mathrm{M}_2(\mathcal{O}_S)$\footnote{Le titre de VI n'est pas écrit.}. Son centre est
$\mathcal{O}_S$, celui de $\mathcal{M}^{\times}$ est $\mathbb{G}_m$ (32), et l'on a la norme
réduite $\det$ et la trace réduite $\mathrm{Tr}$, définies par localisation à partir de
$\mathrm{M}_2$, où elles sont invariantes par automorphismes (33). En marge :
$\det(uv) = \det u \det v$, $\mathrm{Tr}$ est linéaire, $\det$ quadratique,
$\det(u + \lambda) = \lambda^2 + \lambda\,\mathrm{Tr}\,u + \det u$, et la différentielle de
$\det$ en $1$ est $\mathrm{Tr}$. On a des suites exactes
\[
  (34) \qquad 1 \to \mathbb{G}_m \to \mathcal{M}^{\times} \to G \to 1 , \qquad
  (35) \qquad 1 \to \widetilde{G} \to \mathcal{M}^{\times} \xrightarrow{\ \det\ } \mathbb{G}_m \to 1 ,
\]
\[
  (34') \qquad 0 \to \mathcal{O}_S \xrightarrow{\ \lambda \mapsto \lambda\cdot 1\ } \mathcal{M} \to \mathfrak{g} \to 0 , \qquad
  (35') \qquad 0 \to \widetilde{\mathfrak{g}} \to \mathcal{M} \xrightarrow{\ \mathrm{Tr}\ } \mathcal{O}_S \to 0 ,
\]
qui déterminent $G$, $\widetilde{G}$, $\mathfrak{g} = \mathcal{M}/\mathcal{O}_S$ et
$\widetilde{\mathfrak{g}} = \mathcal{M}^{\mathrm{Tr} = 0}$ directement à partir de $\mathcal{M}$ ; les
homomorphismes canoniques sont les composés
\[
  (37)\text{--}(38) \qquad \widetilde{G} \hookrightarrow \mathcal{M}^{\times} \to G , \qquad
  f : \widetilde{\mathfrak{g}} \hookrightarrow \mathcal{M} \to \mathfrak{g} .
\]

\textbf{La forme trace (pages 70 et 71).} La forme
\[
  (39) \qquad \Phi(u, v) = \mathrm{Tr}(uv)
\]
est symétrique et non dégénérée en chaque point — y compris en caractéristique
$2$ —, d'où $\mathcal{M} \simeq \mathcal{M}^{\vee}$ (40). L'orthogonal de $\mathcal{O}_S\cdot 1$ est
$\widetilde{\mathfrak{g}}$, d'où un accouplement parfait
$\widetilde{\mathfrak{g}} \times \mathfrak{g} \to \mathcal{O}_S$ (41) et
\[
  (42) \qquad \widetilde{\mathfrak{g}} \simeq \mathfrak{g}^{\vee} , \qquad \mathfrak{g} \simeq \widetilde{\mathfrak{g}}^{\vee} .
\]
La forme quadratique $Q_{\mathcal{M}}(u) = -\det u$ sur $\mathcal{M}$ est lisse (hyperbolique de
rang $4$), de forme polaire
\[
  (43) \qquad \Phi_1(u, v) = \det u + \det v - \det(u + v) = \mathrm{Tr}(uv) - \mathrm{Tr}\,u\,\mathrm{Tr}\,v ,
\]
elle aussi parfaite. $\Phi_1(1, u) = -\mathrm{Tr}\,u$, donc l'orthogonal de
$\mathcal{O}_S\cdot 1$ pour $\Phi_1$ est encore $\widetilde{\mathfrak{g}}$, et $\Phi$, $\Phi_1$
ont la même restriction $\varphi$ à $\widetilde{\mathfrak{g}}$,
$\varphi(u, v) = \mathrm{Tr}(uv)$. L'accouplement
$\mathfrak{g} \times \widetilde{\mathfrak{g}} \to \mathcal{O}_S$ déduit de $\Phi_1$ est égal à (41),
puisque $\mathrm{Tr}\,v = 0$ pour $v \in \widetilde{\mathfrak{g}}$\footnote{La page dit, sous
une lecture incertaine, « égal au signe près ». Le numéro (42) sert deux fois à
la page 71, et (43) de nouveau à la page 72 ; on garde les numéros de la page.}.

\textbf{La forme $q$ et la forme de Killing (page 72).} Sur $\widetilde{\mathfrak{g}}$,
\[
  (43) \qquad q(u) = -\det u , \qquad
  (44) \qquad q(x + y) = q(x) + q(y) + \varphi(x, y) , \quad \varphi(x, x) = 2q(x) ;
\]
$q$ est lisse, $\varphi$ parfaite sauf en caractéristique $2$. En coordonnées (6),
$q = xy + z^2$ : c'est $q_0$. Et
\[
  (45) \qquad \mathrm{Tr}\bigl(\mathrm{ad}(u)\,\mathrm{ad}(v)\bigr) = 4\varphi(u, v) , \qquad
  \mathrm{Tr}\bigl(\mathrm{ad}(u)^2\bigr) = 8\,q(u) ,
\]
la forme de Killing de $\widetilde{\mathfrak{g}}$ ; c'est (74)--(76) de la page 24,
avec $\Delta = -2$.

\textbf{La base $\omega$ (pages 72 et 73).} Localement $\mathcal{M} \simeq V^{\vee} \otimes V$,
et
\[
  (47) \qquad \det(V^{\vee} \otimes V) \simeq (\det V^{\vee})^{\otimes 2} \otimes (\det V)^{\otimes 2} \simeq \mathcal{O}_S
\]
\footnote{La page écrit $\det V^{\vee} \otimes \det V$, sans les carrés ; le
déterminant d'un produit tensoriel de deux modules de rang $2$ porte l'exposant
$2$. La conclusion, une trivialisation canonique, ne change pas.} ; cet
isomorphisme est stable par $\mathrm{GL}(V)$, donc par $\mathrm{PGL}(V)$, et descend :
d'où une base de $\det\mathcal{M}$, donc, par (35'), une base $\omega$ de
$\det\widetilde{\mathfrak{g}}$. Il y a dans ces identifications une convention de
\emph{signe} à faire — sur $\mathbf{Z}$, un Module inversible trivial a exactement deux
bases, opposées —, mais $\omega^{\otimes 2}$ n'en dépend pas, et
\[
  (48)\text{--}(49) \qquad \delta'(q) = -\,\omega^{\otimes(-2)} , \qquad \delta(\varphi) = -2\,\omega^{\otimes(-2)}
\]
\footnote{La page écrit $\omega^{\otimes 2}$ ; $\delta'(q)$ vit dans
$(\det\widetilde{\mathfrak{g}})^{\otimes(-2)}$, et l'on suit (33) et (82).} : $(\widetilde{\mathfrak{g}},
q, \omega)$ est un module spécial quadratique.

\pagerange{74}{75}
\subsection*{$(E, q, \omega) = (\widetilde{\mathfrak{g}}, q, \omega)$ ; les applications $f$ et $f'$ (pages 74 et 75)}

On pose
\[
  (50) \qquad (E, q, \omega) = (\widetilde{\mathfrak{g}}, q, \omega) .
\]
C'est, en substance, l'isomorphisme $i : E \simeq \widetilde{\mathfrak{g}}$ que la page 31
laissait à faire, pris ici comme définition. L'homomorphisme
$f : \widetilde{\mathfrak{g}} \to \mathfrak{g}$ de (38) n'est autre que l'application
$\widetilde{\mathfrak{g}} \to \widetilde{\mathfrak{g}}^{\vee}$ associée à $\varphi$, compte tenu
de $\widetilde{\mathfrak{g}}^{\vee} \simeq \mathfrak{g}$ (42) : $\langle u, f(v) \rangle =
\mathrm{Tr}(uv) = \varphi(u, v)$. C'est le $f$ de la première suite. Les crochets de
$\widetilde{\mathfrak{g}}$ et de $\mathfrak{g}$ se déduisent alors de $(\varphi, \omega)$ par
les constructions des pages 2 à 9 ; c'est ici que la convention de signe compte,
et l'on prend
\[
  (52) \qquad \omega = \widetilde{X} \wedge \widetilde{H} \wedge \widetilde{Y} ,
\]
avec $\widetilde{X} = E_{12}$, $\widetilde{Y} = E_{21}$,
$\widetilde{H} = \mathrm{diag}(1, -1)$ — en marge : « Sans garantie — il faut
vérifier… ». C'est le bon signe : c'est celui de (64), page 18, et avec lui le
crochet (40 bis) construit à partir de $(\varphi, \omega)$ est exactement le
commutateur de $\mathfrak{sl}_2$, (70), tandis que le crochet (12) sur
$E' = \mathfrak{g}$ est le commutateur de $\mathfrak{gl}_2/\mathcal{O}$\footnote{Vérification
de l'édition, à partir des calculs des pages 18 à 20 : avec $X$, $Y$, $H$ les
classes de $E_{12}$, $E_{21}$, $E_{11}$ dans $\mathfrak{g}$, l'accouplement
$\mathrm{Tr}(uv)$ fait des bases $(\widetilde{Y}, \widetilde{X}, \widetilde{H})$ et
$(X, Y, H)$ des bases duales, comme $(e_1, e_2, e_3)$ et $(e'_1, e'_2, e'_3)$
à la page 18. La page 74 écrit aussi $\omega = -\widetilde{X} \wedge \widetilde{Y} \wedge
\widetilde{H}$ en indiquant sous les lettres $e_1, e_2, e_3$, base de $E_0$ par (6) :
dans les coordonnées (6), qui envoient $e_1$ sur $E_{12}$, (52) est donc
$\omega = -e_1 \wedge e_2 \wedge e_3$, l'opposé de la base choisie pour le représentant
typique de la page 60 ; pour que (6) respecte aussi $\omega$, il faut la composer
avec $z \mapsto -z$, qui préserve $q_0$. L'étiquetage (62) de la page 18, où
$\widetilde{X} = e_2$, donne au contraire $\omega = e_1 \wedge e_2 \wedge e_3$ : les
étiquetages diffèrent, la base $\omega$ de (52) et de (64) est la même. Le dossier 77 (page 89) fixe au contraire $\omega$ par
$\varphi([u, v], w)\,\omega = 2\,u \wedge v \wedge w$, ce qui donne la base opposée : avec
les conventions d'ici, $\varphi([u, v], w)\,\omega = -2\,u \wedge v \wedge w$, puisque
$\varphi([x, y], z)$ vaut $\Delta = -2$ fois le rapport de $x \wedge y \wedge z$ à
$\omega$. C'est exactement l'ambiguïté de signe de la page 73.}.

L'application cofacteur s'explicite :
\[
  (53)\text{--}(55) \qquad f' : \mathfrak{g} \longrightarrow \widetilde{\mathfrak{g}} , \qquad
  f'(\bar{u}) = (\mathrm{Tr}\,u)\cdot 1 - 2u , \qquad ff' = -2\,\mathrm{id}_{\mathfrak{g}} , \quad f'f = -2\,\mathrm{id}_{\widetilde{\mathfrak{g}}} ,
\]
où $\bar{u}$ est l'image dans $\mathfrak{g}$ de $u \in \mathcal{M}$ : comme $\mathrm{Tr}\,1 = 2$,
$u \mapsto \mathrm{Tr}(u) - 2u$ est à valeurs dans $\widetilde{\mathfrak{g}}$ et nulle sur
$\mathcal{O}_S\cdot 1$ ; pour $u \in \widetilde{\mathfrak{g}}$, $f'f(u) = -2u$, et
$ff'(\bar{u}) = -2\bar{u}$. On retrouve $\Delta = -2$ et le signe corrigé de (66) :
$f'(\bar{E}_{12}) = -2E_{12}$, $f'(\bar{E}_{11}) = 1 - 2E_{11} = -\widetilde{H}$.
La page dit $f'$ « s'identifiant à $\Lambda^2 f$ » ; avec le signe (52), c'est
exact.

Enfin le plan et la conique :
\[
  (56) \qquad P = \check{\mathbb{P}}(\widetilde{\mathfrak{g}}) \simeq \mathbb{P}(\mathfrak{g}) , \qquad
  X = \mathrm{Var}(q) \subset P .
\]

\pagerange{76}{77}
\subsection*{La conique et les sous-groupes de Borel (pages 76 et 77)}

La description $X = \mathrm{Var}(q)$ et la description
$X \simeq \underline{\mathrm{Bor}}(G) \simeq \underline{\mathrm{Bor}}(\widetilde{G})$ se raccordent
ainsi.

\textbf{Proposition (page 76).} \emph{Les sous-groupes de Borel $B$ de $G$, $B_1$ de
$\mathcal{M}^{\times}$ et $\widetilde{B}$ de $\widetilde{G}$ se correspondent bijectivement par
image inverse. Un sous-groupe de Borel est connu par son algèbre de Lie : celle
$\mathfrak{L}$ de $B$, sous-fibré de corang $1$ de $\mathfrak{g}$, ou celle
$\mathfrak{L}_1$ de $B_1$, sous-fibré de corang $1$ de $\mathcal{M}$, image inverse de
$\mathfrak{L}$, qui est une sous-algèbre ; et, si $2$ est inversible, celle
$\widetilde{\mathfrak{L}} = \mathfrak{L}_1 \cap \widetilde{\mathfrak{g}}$ de $\widetilde{B}$, puisque
$f$ est alors un isomorphisme. Explicitement}
\[
  (57) \qquad B_1 = \mathfrak{L}_1^{\times} = \mathbb{V}(\mathfrak{L}_1) \cap \mathcal{M}^{\times} , \qquad
  B = B_1/\mathbb{G}_m , \qquad \widetilde{B} = B_1 \cap \widetilde{G} .
\]
\footnote{La page écrit $\widetilde{B} = \check{\mathbb{V}}(\widetilde{\mathfrak{L}}) \cap
\mathcal{M}^{\times}$ et $B = \widetilde{B}/\mathbb{G}_m$ : le groupe qu'elle a en vue est $B_1$,
le groupe des unités de la sous-algèbre $\mathfrak{L}_1$ de rang $3$ (localement,
les matrices triangulaires supérieures). Les inversibles du module
$\widetilde{\mathfrak{L}}$, de rang $2$, ne forment pas un groupe. La lettre lue
$\widetilde{B}$ porte peut-être une barre sur la page.}

\textbf{Critère (page 77).} Si $2$ est inversible, un sous-fibré
$\mathfrak{L} \subset \widetilde{\mathfrak{g}}$ de corang $1$ est l'algèbre de Lie d'un Borel
si et seulement si son orthogonal $\mathfrak{L}^{\perp}$ (pour $\varphi$), de rang $1$,
est isotrope pour $q$ : localement, le Borel triangulaire supérieur a pour
algèbre de Lie les $\begin{pmatrix} a & b \\ 0 & -a \end{pmatrix}$, d'orthogonal
$\mathcal{O} E_{12}$, d'élément de carré nul. La droite isotrope $\mathfrak{L}^{\perp}$ est
le point de la conique qui correspond au Borel : c'est le raccord cherché. La
page ajoute que cela caractérise $q$ à un scalaire inversible près, et pose en
marge la question : « suffit-il, pour $\mathfrak{L}$, d'être une sous-algèbre ? »,
qu'elle laisse ouverte\footnote{La fin de la parenthèse de la page 77, chargée
d'insertions, se lit mal ; une seconde note marginale parle d'éléments de « carré
nul » de $\mathcal{M}$. L'hypothèse « $2$ inversible » est de l'édition : c'est celle
sous laquelle $\varphi$ est parfaite sur $\widetilde{\mathfrak{g}}$, et la page passe
aussitôt au cas de caractéristique $2$.}.

\pagerange{77}{78}
\subsection*{Caractéristique $2$ : le Frobenius (pages 77 et 78)}

Supposons $2 = 0$ sur $S$. Alors $\mathrm{Tr}\,1 = 2 = 0$, donc
\[
  (58) \qquad \mathcal{O}_S\cdot 1 \subset \widetilde{\mathfrak{g}} ,
\]
et c'est le noyau de $\varphi$, c'est-à-dire de $f$. Le morphisme
\[
  (59) \qquad X \longrightarrow \mathbb{P}(\widetilde{\mathfrak{g}}) ,
\]
qui envoie un Borel sur son algèbre de Lie — c'est-à-dire une droite isotrope
$L \subset \widetilde{\mathfrak{g}}$ sur son orthogonal pour $\varphi$ —, prend ses
valeurs dans les hyperplans qui contiennent $\mathcal{O}_S\cdot 1$, et se factorise donc
par un morphisme de fibrés en droites projectives
\[
  (60) \qquad X \longrightarrow \mathbb{P}(\widetilde{\mathfrak{g}}/\mathcal{O}_S) .
\]
Ce n'est pas un monomorphisme : s'il l'était, ce serait un isomorphisme, et $X$
serait banal dès qu'on est en caractéristique $2$, alors qu'il existe des fibrés
en droites projectives non banals en caractéristique $2$. En fait il y a un
isomorphisme canonique
\[
  (61) \qquad \mathbb{P}(\widetilde{\mathfrak{g}}/\mathcal{O}_S) \simeq X^{(2)}
\]
avec le « frobénisé » de $X$, et (60) est le morphisme de Frobenius relatif
$X \to X^{(2)}$, épimorphisme radiciel, fini et localement libre de rang
$2^d = 2$, nullement un monomorphisme\footnote{Vérification de l'édition dans le
cas standard : la conique $xy = z^2$ (caractéristique $2$) est paramétrée par
$(s : t) \mapsto (s^2, t^2, st)$ ; l'orthogonal de la droite isotrope
$\mathcal{O}(s^2\widetilde{X} + t^2\widetilde{Y} + st\,\widetilde{H})$ est l'hyperplan
$t^2 x + s^2 y = 0$, qui contient $\widetilde{H} = 1$ et ne dépend que de
$(s^2 : t^2)$.}. La page, et le dossier, s'arrêtent là.

\end{document}
